%%%%%%%%%%%%%%%%%%%%%%%%%%%%%%%%%%%%%%%%%%%%%%%%%%%%%%%%%%%%
% 13.3 MATHEMATICAL PHYSICS
% The Composition of Advanced Mathematics
%%%%%%%%%%%%%%%%%%%%%%%%%%%%%%%%%%%%%%%%%%%%%%%%%%%%%%%%%%%%

\section{Mathematical Physics}
\label{sec:mathematical-physics}

\prerequisite{
Calculus, multivariable calculus, vector calculus, differential equations,
partial differential equations, linear algebra, complex analysis,
Fourier analysis, variational methods, probability, and mathematical modeling.
}

\learningobjectives{
By the end of this section, the reader should be able to:
\begin{itemize}
    \item explain the role of mathematical models in physics;
    \item distinguish kinematic and dynamical descriptions;
    \item work with position, velocity, acceleration, momentum, and force;
    \item formulate Newtonian equations of motion;
    \item understand conservative force fields;
    \item define potential energy;
    \item derive conservation of mechanical energy;
    \item understand work and power;
    \item recognize central-force motion;
    \item understand angular momentum conservation;
    \item recognize harmonic oscillators;
    \item distinguish free, damped, and forced oscillations;
    \item understand resonance;
    \item recognize coupled oscillators and normal modes;
    \item formulate Lagrangian mechanics;
    \item derive Euler--Lagrange equations;
    \item understand generalized coordinates;
    \item recognize generalized momentum;
    \item understand Hamiltonian mechanics;
    \item derive Hamilton's equations;
    \item recognize phase space;
    \item understand Poisson brackets conceptually;
    \item recognize canonical transformations conceptually;
    \item understand action principles;
    \item formulate basic field equations;
    \item recognize scalar and vector fields;
    \item understand gradient, divergence, curl, and Laplacian operators;
    \item recognize conservation laws in differential form;
    \item formulate wave equations;
    \item understand wave speed;
    \item recognize standing waves;
    \item understand normal modes and eigenvalue problems;
    \item formulate the heat equation;
    \item understand diffusion and thermal transport;
    \item formulate Laplace and Poisson equations;
    \item understand boundary-value problems;
    \item recognize Green's functions conceptually;
    \item understand Fourier methods in physical systems;
    \item recognize electromagnetic field equations;
    \item understand Maxwell's equations conceptually;
    \item derive the electromagnetic wave equation conceptually;
    \item understand electrostatic potential;
    \item recognize gravitational potentials;
    \item understand inverse-square fields;
    \item recognize continuum field descriptions;
    \item understand stress and strain conceptually;
    \item recognize conservation of mass, momentum, and energy;
    \item understand variational formulations of PDEs;
    \item recognize weak formulations conceptually;
    \item understand classical relativity conceptually;
    \item recognize Lorentz transformations;
    \item understand relativistic energy and momentum;
    \item recognize spacetime intervals;
    \item understand basic general-relativistic geometric ideas conceptually;
    \item recognize quantum states and wavefunctions;
    \item understand normalization of wavefunctions;
    \item formulate the Schrödinger equation;
    \item recognize operators and observables;
    \item understand eigenvalues as possible measurement outcomes;
    \item recognize expectation values;
    \item understand uncertainty relations conceptually;
    \item recognize stationary states;
    \item solve basic infinite-well problems;
    \item recognize harmonic-oscillator quantization conceptually;
    \item understand superposition;
    \item recognize Hilbert-space structure in quantum mechanics;
    \item understand symmetry in physical mathematics;
    \item recognize Noether's theorem conceptually;
    \item connect symmetry with conservation laws;
    \item understand dimensional analysis and nondimensionalization in physics;
    \item recognize scaling and similarity;
    \item understand perturbation methods conceptually;
    \item recognize asymptotic approximation;
    \item understand numerical simulation in physics;
    \item distinguish mathematical model error from numerical error;
    \item connect mathematical physics with mechanics, continuum mechanics,
          fluid dynamics, control, differential equations, geometry,
          probability, and computation.
\end{itemize}
}

%%%%%%%%%%%%%%%%%%%%%%%%%%%%%%%%%%%%%%%%%%%%%%%%%%%%%%%%%%%%
\subsection{What Is Mathematical Physics?}
\label{subsec:what-is-mathematical-physics}
%%%%%%%%%%%%%%%%%%%%%%%%%%%%%%%%%%%%%%%%%%%%%%%%%%%%%%%%%%%%

Mathematical physics develops mathematical structures for describing
physical laws and physical systems.

The subject includes mathematics arising from:

\[
\boxed{
\text{mechanics}
}
\]

\[
\boxed{
\text{electromagnetism}
}
\]

\[
\boxed{
\text{waves}
}
\]

\[
\boxed{
\text{heat and diffusion}
}
\]

\[
\boxed{
\text{relativity}
}
\]

and

\[
\boxed{
\text{quantum mechanics}.
}
\]

The same mathematical structures frequently appear in several different
physical contexts.

%%%%%%%%%%%%%%%%%%%%%%%%%%%%%%%%%%%%%%%%%%%%%%%%%%%%%%%%%%%%
\subsection{Physical Quantities and Units}
\label{subsec:physical-quantities-units}
%%%%%%%%%%%%%%%%%%%%%%%%%%%%%%%%%%%%%%%%%%%%%%%%%%%%%%%%%%%%

A physical variable has both a numerical value and a dimension.

Common dimensions include

\[
\boxed{
L
=
\text{length},
}
\]

\[
\boxed{
T
=
\text{time},
}
\]

and

\[
\boxed{
M
=
\text{mass}.
}
\]

For example,

\[
[v]
=
LT^{-1},
\]

and

\[
[a]
=
LT^{-2}.
\]

A physically meaningful equation must be dimensionally consistent.

%%%%%%%%%%%%%%%%%%%%%%%%%%%%%%%%%%%%%%%%%%%%%%%%%%%%%%%%%%%%
\subsection{Position, Velocity, and Acceleration}
\label{subsec:position-velocity-acceleration}
%%%%%%%%%%%%%%%%%%%%%%%%%%%%%%%%%%%%%%%%%%%%%%%%%%%%%%%%%%%%

Let

\[
\mathbf{r}(t)
\]

be the position of a particle.

Velocity is

\[
\boxed{
\mathbf{v}(t)
=
\frac{d\mathbf{r}}{dt}.
}
\]

Acceleration is

\[
\boxed{
\mathbf{a}(t)
=
\frac{d\mathbf{v}}{dt}
=
\frac{d^2\mathbf{r}}{dt^2}.
}
\]

These are purely kinematic definitions.

%%%%%%%%%%%%%%%%%%%%%%%%%%%%%%%%%%%%%%%%%%%%%%%%%%%%%%%%%%%%
\subsection{Newton's Second Law}
\label{subsec:newtons-second-law}
%%%%%%%%%%%%%%%%%%%%%%%%%%%%%%%%%%%%%%%%%%%%%%%%%%%%%%%%%%%%

For constant mass \(m\), Newton's second law is

\[
\boxed{
m\mathbf{a}
=
\mathbf{F}.
}
\]

Thus the motion satisfies

\[
\boxed{
m\frac{d^2\mathbf{r}}{dt^2}
=
\mathbf{F}
\left(
\mathbf{r},
\mathbf{v},
t
\right).
}
\]

This converts a physical force law into a differential equation.

%%%%%%%%%%%%%%%%%%%%%%%%%%%%%%%%%%%%%%%%%%%%%%%%%%%%%%%%%%%%
\subsection{Momentum}
\label{subsec:classical-momentum}
%%%%%%%%%%%%%%%%%%%%%%%%%%%%%%%%%%%%%%%%%%%%%%%%%%%%%%%%%%%%

Linear momentum is

\[
\boxed{
\mathbf{p}
=
m\mathbf{v}.
}
\]

More generally, Newton's law can be written

\[
\boxed{
\frac{d\mathbf{p}}{dt}
=
\mathbf{F}.
}
\]

If total force vanishes,

\[
\mathbf{F}=0,
\]

then

\[
\boxed{
\mathbf{p}
=
\text{constant}.
}
\]

%%%%%%%%%%%%%%%%%%%%%%%%%%%%%%%%%%%%%%%%%%%%%%%%%%%%%%%%%%%%
\subsection{Worked Example: Constant Force}
\label{subsec:worked-constant-force}
%%%%%%%%%%%%%%%%%%%%%%%%%%%%%%%%%%%%%%%%%%%%%%%%%%%%%%%%%%%%

\begin{workedexamplebox}{Particle Under Constant Force}

Suppose a one-dimensional particle of mass \(m\) experiences constant force
\(F_0\).

Newton's law gives

\[
m\frac{d^2x}{dt^2}
=
F_0.
\]

Therefore,

\[
\frac{d^2x}{dt^2}
=
\frac{F_0}{m}.
\]

Integrating,

\[
v(t)
=
v_0
+
\frac{F_0}{m}t.
\]

Integrating again,

\[
x(t)
=
x_0
+
v_0t
+
\frac{F_0}{2m}t^2.
\]

\begin{answerbox}
\[
\boxed{
x(t)
=
x_0
+
v_0t
+
\frac{F_0}{2m}t^2.
}
\]
\end{answerbox}

\end{workedexamplebox}

%%%%%%%%%%%%%%%%%%%%%%%%%%%%%%%%%%%%%%%%%%%%%%%%%%%%%%%%%%%%
\subsection{Work}
\label{subsec:physical-work}
%%%%%%%%%%%%%%%%%%%%%%%%%%%%%%%%%%%%%%%%%%%%%%%%%%%%%%%%%%%%

The work done by a force along path \(C\) is

\[
\boxed{
W
=
\int_C
\mathbf{F}\cdot d\mathbf{r}.
}
\]

In one dimension,

\[
\boxed{
W
=
\int_{x_1}^{x_2}
F(x)\,dx.
}
\]

Work measures energy transferred through force acting over displacement.

%%%%%%%%%%%%%%%%%%%%%%%%%%%%%%%%%%%%%%%%%%%%%%%%%%%%%%%%%%%%
\subsection{Kinetic Energy}
\label{subsec:kinetic-energy}
%%%%%%%%%%%%%%%%%%%%%%%%%%%%%%%%%%%%%%%%%%%%%%%%%%%%%%%%%%%%

For a particle of mass \(m\),

\[
\boxed{
T
=
\frac12mv^2.
}
\]

Here \(T\) conventionally denotes kinetic energy in analytical mechanics.

The work-energy theorem states

\[
\boxed{
W_{\text{net}}
=
\Delta T.
}
\]

%%%%%%%%%%%%%%%%%%%%%%%%%%%%%%%%%%%%%%%%%%%%%%%%%%%%%%%%%%%%
\subsection{Derivation of the Work-Energy Theorem}
\label{subsec:work-energy-derivation}
%%%%%%%%%%%%%%%%%%%%%%%%%%%%%%%%%%%%%%%%%%%%%%%%%%%%%%%%%%%%

In one dimension,

\[
F
=
m\frac{dv}{dt}.
\]

Since

\[
\frac{dx}{dt}
=
v,
\]

we obtain

\[
F\,dx
=
m\frac{dv}{dt}v\,dt
=
mv\,dv.
\]

Integrating,

\[
W
=
\int_{v_1}^{v_2}
mv\,dv.
\]

Therefore,

\[
\boxed{
W
=
\frac12mv_2^2
-
\frac12mv_1^2.
}
\]

%%%%%%%%%%%%%%%%%%%%%%%%%%%%%%%%%%%%%%%%%%%%%%%%%%%%%%%%%%%%
\subsection{Conservative Forces}
\label{subsec:conservative-forces}
%%%%%%%%%%%%%%%%%%%%%%%%%%%%%%%%%%%%%%%%%%%%%%%%%%%%%%%%%%%%

A force is conservative when it can be written as the negative gradient of
a scalar potential \(V\):

\[
\boxed{
\mathbf{F}
=
-\nabla V.
}
\]

In one dimension,

\[
\boxed{
F(x)
=
-\frac{dV}{dx}.
}
\]

The function \(V\) is potential energy.

%%%%%%%%%%%%%%%%%%%%%%%%%%%%%%%%%%%%%%%%%%%%%%%%%%%%%%%%%%%%
\subsection{Mechanical Energy}
\label{subsec:mechanical-energy}
%%%%%%%%%%%%%%%%%%%%%%%%%%%%%%%%%%%%%%%%%%%%%%%%%%%%%%%%%%%%

For a conservative system,

\[
\boxed{
E
=
T+V.
}
\]

In one dimension,

\[
E
=
\frac12m\dot{x}^2
+
V(x).
\]

When \(V\) has no explicit time dependence and no nonconservative forces are
present,

\[
\boxed{
\frac{dE}{dt}=0.
}
\]

Thus mechanical energy is conserved.

%%%%%%%%%%%%%%%%%%%%%%%%%%%%%%%%%%%%%%%%%%%%%%%%%%%%%%%%%%%%
\subsection{Worked Example: Gravitational Potential Near Earth}
\label{subsec:worked-gravitational-potential}
%%%%%%%%%%%%%%%%%%%%%%%%%%%%%%%%%%%%%%%%%%%%%%%%%%%%%%%%%%%%

\begin{workedexamplebox}{Vertical Motion}

Near Earth's surface, approximate gravitational force by

\[
F=-mg.
\]

Since

\[
F
=
-\frac{dV}{dy},
\]

we have

\[
-mg
=
-\frac{dV}{dy}.
\]

Thus

\[
\frac{dV}{dy}
=
mg.
\]

Integrating,

\[
V(y)=mgy+C.
\]

Choosing \(C=0\),

\begin{answerbox}
\[
\boxed{
V(y)=mgy.
}
\]
\end{answerbox}

\end{workedexamplebox}

%%%%%%%%%%%%%%%%%%%%%%%%%%%%%%%%%%%%%%%%%%%%%%%%%%%%%%%%%%%%
\subsection{Power}
\label{subsec:physical-power}
%%%%%%%%%%%%%%%%%%%%%%%%%%%%%%%%%%%%%%%%%%%%%%%%%%%%%%%%%%%%

Power is the rate at which work is performed:

\[
\boxed{
P
=
\frac{dW}{dt}.
}
\]

Since

\[
dW
=
\mathbf{F}\cdot d\mathbf{r},
\]

we obtain

\[
\boxed{
P
=
\mathbf{F}\cdot\mathbf{v}.
}
\]

%%%%%%%%%%%%%%%%%%%%%%%%%%%%%%%%%%%%%%%%%%%%%%%%%%%%%%%%%%%%
\subsection{Angular Momentum}
\label{subsec:angular-momentum}
%%%%%%%%%%%%%%%%%%%%%%%%%%%%%%%%%%%%%%%%%%%%%%%%%%%%%%%%%%%%

Angular momentum about the origin is

\[
\boxed{
\mathbf{L}
=
\mathbf{r}\times\mathbf{p}.
}
\]

Its time derivative is

\[
\boxed{
\frac{d\mathbf{L}}{dt}
=
\boldsymbol{\tau},
}
\]

where

\[
\boldsymbol{\tau}
=
\mathbf{r}\times\mathbf{F}
\]

is torque.

If

\[
\boldsymbol{\tau}=0,
\]

then angular momentum is conserved.

%%%%%%%%%%%%%%%%%%%%%%%%%%%%%%%%%%%%%%%%%%%%%%%%%%%%%%%%%%%%
\subsection{Central Forces}
\label{subsec:central-forces}
%%%%%%%%%%%%%%%%%%%%%%%%%%%%%%%%%%%%%%%%%%%%%%%%%%%%%%%%%%%%

A central force points along the radial direction:

\[
\boxed{
\mathbf{F}
=
F(r)\hat{\mathbf{r}}.
}
\]

Since

\[
\mathbf{r}\times\mathbf{F}=0,
\]

the torque vanishes.

Therefore,

\[
\boxed{
\mathbf{L}
=
\text{constant}.
}
\]

The motion is confined to a plane perpendicular to \(\mathbf{L}\).

%%%%%%%%%%%%%%%%%%%%%%%%%%%%%%%%%%%%%%%%%%%%%%%%%%%%%%%%%%%%
\subsection{Inverse-Square Forces}
\label{subsec:inverse-square-forces}
%%%%%%%%%%%%%%%%%%%%%%%%%%%%%%%%%%%%%%%%%%%%%%%%%%%%%%%%%%%%

Gravitational and electrostatic forces provide important inverse-square
examples.

A radial inverse-square force has magnitude

\[
\boxed{
F(r)
=
\frac{k}{r^2}.
}
\]

For Newtonian gravity between masses \(M\) and \(m\),

\[
\boxed{
\mathbf{F}
=
-\frac{GMm}{r^2}
\hat{\mathbf{r}}.
}
\]

%%%%%%%%%%%%%%%%%%%%%%%%%%%%%%%%%%%%%%%%%%%%%%%%%%%%%%%%%%%%
\subsection{Gravitational Potential}
\label{subsec:newtonian-gravitational-potential}
%%%%%%%%%%%%%%%%%%%%%%%%%%%%%%%%%%%%%%%%%%%%%%%%%%%%%%%%%%%%

For gravitational force

\[
\mathbf{F}
=
-\frac{GMm}{r^2}\hat{\mathbf{r}},
\]

the potential energy is

\[
\boxed{
V(r)
=
-\frac{GMm}{r},
}
\]

choosing

\[
V(\infty)=0.
\]

Indeed,

\[
-\frac{dV}{dr}
=
-\frac{GMm}{r^2}.
\]

%%%%%%%%%%%%%%%%%%%%%%%%%%%%%%%%%%%%%%%%%%%%%%%%%%%%%%%%%%%%
\subsection{Simple Harmonic Oscillator}
\label{subsec:simple-harmonic-oscillator}
%%%%%%%%%%%%%%%%%%%%%%%%%%%%%%%%%%%%%%%%%%%%%%%%%%%%%%%%%%%%

Hooke's law gives

\[
\boxed{
F=-kx.
}
\]

Newton's law becomes

\[
m\ddot{x}
=
-kx.
\]

Therefore,

\[
\boxed{
\ddot{x}
+
\omega_0^2x
=
0,
}
\]

where

\[
\boxed{
\omega_0
=
\sqrt{\frac{k}{m}}.
}
\]

%%%%%%%%%%%%%%%%%%%%%%%%%%%%%%%%%%%%%%%%%%%%%%%%%%%%%%%%%%%%
\subsection{Oscillator Solution}
\label{subsec:harmonic-oscillator-solution}
%%%%%%%%%%%%%%%%%%%%%%%%%%%%%%%%%%%%%%%%%%%%%%%%%%%%%%%%%%%%

The general solution is

\[
\boxed{
x(t)
=
A\cos(\omega_0t)
+
B\sin(\omega_0t).
}
\]

Equivalently,

\[
\boxed{
x(t)
=
C\cos(\omega_0t-\phi).
}
\]

The period is

\[
\boxed{
T
=
\frac{2\pi}{\omega_0}.
}
\]

%%%%%%%%%%%%%%%%%%%%%%%%%%%%%%%%%%%%%%%%%%%%%%%%%%%%%%%%%%%%
\subsection{Oscillator Energy}
\label{subsec:harmonic-oscillator-energy}
%%%%%%%%%%%%%%%%%%%%%%%%%%%%%%%%%%%%%%%%%%%%%%%%%%%%%%%%%%%%

The potential energy is

\[
\boxed{
V(x)
=
\frac12kx^2.
}
\]

The total energy is

\[
\boxed{
E
=
\frac12m\dot{x}^2
+
\frac12kx^2.
}
\]

In the undamped oscillator,

\[
\boxed{
E
=
\text{constant}.
}
\]

%%%%%%%%%%%%%%%%%%%%%%%%%%%%%%%%%%%%%%%%%%%%%%%%%%%%%%%%%%%%
\subsection{Worked Example: Oscillator Frequency}
\label{subsec:worked-oscillator-frequency}
%%%%%%%%%%%%%%%%%%%%%%%%%%%%%%%%%%%%%%%%%%%%%%%%%%%%%%%%%%%%

\begin{workedexamplebox}{Mass-Spring Frequency}

Suppose

\[
m=2,
\qquad
k=18.
\]

Then

\[
\omega_0
=
\sqrt{\frac{18}{2}}
=
3.
\]

The period is

\[
T
=
\frac{2\pi}{3}.
\]

\begin{answerbox}
\[
\boxed{
\omega_0=3,
\qquad
T=\frac{2\pi}{3}.
}
\]
\end{answerbox}

\end{workedexamplebox}

%%%%%%%%%%%%%%%%%%%%%%%%%%%%%%%%%%%%%%%%%%%%%%%%%%%%%%%%%%%%
\subsection{Damped Oscillator}
\label{subsec:damped-oscillator}
%%%%%%%%%%%%%%%%%%%%%%%%%%%%%%%%%%%%%%%%%%%%%%%%%%%%%%%%%%%%

Linear damping gives

\[
\boxed{
m\ddot{x}
+
c\dot{x}
+
kx
=
0.
}
\]

The characteristic equation is

\[
\boxed{
mr^2+cr+k=0.
}
\]

The behavior depends on the discriminant

\[
c^2-4mk.
\]

%%%%%%%%%%%%%%%%%%%%%%%%%%%%%%%%%%%%%%%%%%%%%%%%%%%%%%%%%%%%
\subsection{Damping Regimes}
\label{subsec:damping-regimes}
%%%%%%%%%%%%%%%%%%%%%%%%%%%%%%%%%%%%%%%%%%%%%%%%%%%%%%%%%%%%

If

\[
\boxed{
c^2<4mk,
}
\]

the system is underdamped.

If

\[
\boxed{
c^2=4mk,
}
\]

the system is critically damped.

If

\[
\boxed{
c^2>4mk,
}
\]

the system is overdamped.

These regimes describe qualitatively different returns toward equilibrium.

%%%%%%%%%%%%%%%%%%%%%%%%%%%%%%%%%%%%%%%%%%%%%%%%%%%%%%%%%%%%
\subsection{Forced Oscillator}
\label{subsec:forced-oscillator}
%%%%%%%%%%%%%%%%%%%%%%%%%%%%%%%%%%%%%%%%%%%%%%%%%%%%%%%%%%%%

An external periodic force gives

\[
\boxed{
m\ddot{x}
+
c\dot{x}
+
kx
=
F_0\cos(\omega t).
}
\]

After transients decay, the steady-state response oscillates at driving
frequency \(\omega\).

Its amplitude depends strongly on the relationship between \(\omega\) and
the system's natural frequency.

%%%%%%%%%%%%%%%%%%%%%%%%%%%%%%%%%%%%%%%%%%%%%%%%%%%%%%%%%%%%
\subsection{Resonance}
\label{subsec:physical-resonance}
%%%%%%%%%%%%%%%%%%%%%%%%%%%%%%%%%%%%%%%%%%%%%%%%%%%%%%%%%%%%

Resonance occurs when a periodic driving force produces a particularly
large response.

For weak damping, the strongest response occurs near

\[
\boxed{
\omega
\approx
\omega_0.
}
\]

Resonance appears in:

\[
\boxed{
\text{mechanical systems},
}
\]

\[
\boxed{
\text{electrical circuits},
}
\]

\[
\boxed{
\text{acoustics},
}
\]

and

\[
\boxed{
\text{wave systems}.
}
\]

%%%%%%%%%%%%%%%%%%%%%%%%%%%%%%%%%%%%%%%%%%%%%%%%%%%%%%%%%%%%
\subsection{Coupled Oscillators}
\label{subsec:coupled-oscillators}
%%%%%%%%%%%%%%%%%%%%%%%%%%%%%%%%%%%%%%%%%%%%%%%%%%%%%%%%%%%%

For several interacting oscillators,

\[
\boxed{
M\ddot{\mathbf{x}}
+
K\mathbf{x}
=
0,
}
\]

where \(M\) is the mass matrix and \(K\) is the stiffness matrix.

Seek solutions of the form

\[
\boxed{
\mathbf{x}(t)
=
\mathbf{v}e^{i\omega t}.
}
\]

Substitution gives

\[
\boxed{
(K-\omega^2M)\mathbf{v}
=
0.
}
\]

%%%%%%%%%%%%%%%%%%%%%%%%%%%%%%%%%%%%%%%%%%%%%%%%%%%%%%%%%%%%
\subsection{Normal Modes}
\label{subsec:normal-modes}
%%%%%%%%%%%%%%%%%%%%%%%%%%%%%%%%%%%%%%%%%%%%%%%%%%%%%%%%%%%%

A nonzero normal-mode vector \(\mathbf{v}\) exists when

\[
\boxed{
\det(K-\omega^2M)=0.
}
\]

Thus normal frequencies are obtained from a generalized eigenvalue problem.

Each eigenvector gives a characteristic pattern of collective oscillation.

%%%%%%%%%%%%%%%%%%%%%%%%%%%%%%%%%%%%%%%%%%%%%%%%%%%%%%%%%%%%
\subsection{Lagrangian Mechanics}
\label{subsec:lagrangian-mechanics}
%%%%%%%%%%%%%%%%%%%%%%%%%%%%%%%%%%%%%%%%%%%%%%%%%%%%%%%%%%%%

Lagrangian mechanics describes dynamics using generalized coordinates

\[
q_1,\ldots,q_n.
\]

The Lagrangian is often

\[
\boxed{
L(q,\dot q,t)
=
T-V,
}
\]

where \(T\) is kinetic energy and \(V\) is potential energy.

The equations of motion follow from a variational principle.

%%%%%%%%%%%%%%%%%%%%%%%%%%%%%%%%%%%%%%%%%%%%%%%%%%%%%%%%%%%%
\subsection{Action}
\label{subsec:physical-action}
%%%%%%%%%%%%%%%%%%%%%%%%%%%%%%%%%%%%%%%%%%%%%%%%%%%%%%%%%%%%

The action is

\[
\boxed{
S[q]
=
\int_{t_1}^{t_2}
L(q,\dot q,t)\,dt.
}
\]

The physical path makes the action stationary under appropriate variations:

\[
\boxed{
\delta S=0.
}
\]

This principle yields the Euler--Lagrange equations.

%%%%%%%%%%%%%%%%%%%%%%%%%%%%%%%%%%%%%%%%%%%%%%%%%%%%%%%%%%%%
\subsection{Euler--Lagrange Equations}
\label{subsec:physical-euler-lagrange}
%%%%%%%%%%%%%%%%%%%%%%%%%%%%%%%%%%%%%%%%%%%%%%%%%%%%%%%%%%%%

For generalized coordinate \(q_i\),

\[
\boxed{
\frac{d}{dt}
\left(
\frac{\partial L}{\partial\dot q_i}
\right)
-
\frac{\partial L}{\partial q_i}
=
0.
}
\]

These equations reproduce Newtonian mechanics while allowing more flexible
coordinates and constraints.

%%%%%%%%%%%%%%%%%%%%%%%%%%%%%%%%%%%%%%%%%%%%%%%%%%%%%%%%%%%%
\subsection{Worked Example: Harmonic Oscillator From a Lagrangian}
\label{subsec:worked-lagrangian-oscillator}
%%%%%%%%%%%%%%%%%%%%%%%%%%%%%%%%%%%%%%%%%%%%%%%%%%%%%%%%%%%%

\begin{workedexamplebox}{Recovering the Oscillator Equation}

Take

\[
L
=
\frac12m\dot{x}^2
-
\frac12kx^2.
\]

Then

\[
\frac{\partial L}{\partial\dot{x}}
=
m\dot{x},
\]

so

\[
\frac{d}{dt}
\left(
\frac{\partial L}{\partial\dot{x}}
\right)
=
m\ddot{x}.
\]

Also,

\[
\frac{\partial L}{\partial x}
=
-kx.
\]

Therefore the Euler--Lagrange equation gives

\[
m\ddot{x}+kx=0.
\]

\begin{answerbox}
\[
\boxed{
m\ddot{x}+kx=0.
}
\]
\end{answerbox}

\end{workedexamplebox}

%%%%%%%%%%%%%%%%%%%%%%%%%%%%%%%%%%%%%%%%%%%%%%%%%%%%%%%%%%%%
\subsection{Generalized Momentum}
\label{subsec:generalized-momentum}
%%%%%%%%%%%%%%%%%%%%%%%%%%%%%%%%%%%%%%%%%%%%%%%%%%%%%%%%%%%%

The momentum conjugate to \(q_i\) is

\[
\boxed{
p_i
=
\frac{\partial L}{\partial\dot q_i}.
}
\]

For a standard Cartesian particle,

\[
L
=
\frac12m\dot{\mathbf{r}}^2
-
V(\mathbf{r}),
\]

this reduces to ordinary momentum

\[
\boxed{
\mathbf{p}
=
m\dot{\mathbf{r}}.
}
\]

%%%%%%%%%%%%%%%%%%%%%%%%%%%%%%%%%%%%%%%%%%%%%%%%%%%%%%%%%%%%
\subsection{Cyclic Coordinates}
\label{subsec:cyclic-coordinates}
%%%%%%%%%%%%%%%%%%%%%%%%%%%%%%%%%%%%%%%%%%%%%%%%%%%%%%%%%%%%

If the Lagrangian does not depend explicitly on coordinate \(q_i\), then

\[
\frac{\partial L}{\partial q_i}=0.
\]

Euler--Lagrange gives

\[
\frac{d}{dt}
\left(
\frac{\partial L}{\partial\dot q_i}
\right)
=
0.
\]

Therefore,

\[
\boxed{
p_i
=
\text{constant}.
}
\]

A symmetry in a coordinate produces a conserved momentum.

%%%%%%%%%%%%%%%%%%%%%%%%%%%%%%%%%%%%%%%%%%%%%%%%%%%%%%%%%%%%
\subsection{Hamiltonian Mechanics}
\label{subsec:hamiltonian-mechanics}
%%%%%%%%%%%%%%%%%%%%%%%%%%%%%%%%%%%%%%%%%%%%%%%%%%%%%%%%%%%%

The Hamiltonian is defined by the Legendre transform

\[
\boxed{
H(q,p,t)
=
\sum_i p_i\dot q_i
-
L.
}
\]

For many standard conservative systems,

\[
\boxed{
H=T+V.
}
\]

The Hamiltonian often represents total energy.

%%%%%%%%%%%%%%%%%%%%%%%%%%%%%%%%%%%%%%%%%%%%%%%%%%%%%%%%%%%%
\subsection{Hamilton's Equations}
\label{subsec:hamiltons-equations}
%%%%%%%%%%%%%%%%%%%%%%%%%%%%%%%%%%%%%%%%%%%%%%%%%%%%%%%%%%%%

Hamilton's equations are

\[
\boxed{
\dot q_i
=
\frac{\partial H}{\partial p_i},
}
\]

\[
\boxed{
\dot p_i
=
-\frac{\partial H}{\partial q_i}.
}
\]

Thus second-order Newtonian dynamics become a first-order system on phase
space.

%%%%%%%%%%%%%%%%%%%%%%%%%%%%%%%%%%%%%%%%%%%%%%%%%%%%%%%%%%%%
\subsection{Worked Example: Hamiltonian Oscillator}
\label{subsec:worked-hamiltonian-oscillator}
%%%%%%%%%%%%%%%%%%%%%%%%%%%%%%%%%%%%%%%%%%%%%%%%%%%%%%%%%%%%

\begin{workedexamplebox}{Hamilton's Equations}

For a harmonic oscillator,

\[
H(x,p)
=
\frac{p^2}{2m}
+
\frac12kx^2.
\]

Then

\[
\dot{x}
=
\frac{\partial H}{\partial p}
=
\frac{p}{m},
\]

and

\[
\dot{p}
=
-\frac{\partial H}{\partial x}
=
-kx.
\]

Differentiate the first equation:

\[
\ddot{x}
=
\frac{\dot p}{m}.
\]

Therefore,

\[
\begin{answerbox}
\[
\boxed{
m\ddot{x}+kx=0.
}
\]
\end{answerbox}

\end{workedexamplebox}

%%%%%%%%%%%%%%%%%%%%%%%%%%%%%%%%%%%%%%%%%%%%%%%%%%%%%%%%%%%%
\subsection{Phase Space}
\label{subsec:physical-phase-space}
%%%%%%%%%%%%%%%%%%%%%%%%%%%%%%%%%%%%%%%%%%%%%%%%%%%%%%%%%%%%

Phase space consists of generalized coordinates and momenta:

\[
\boxed{
(q_1,\ldots,q_n,p_1,\ldots,p_n).
}
\]

A dynamical state is represented by one point in this \(2n\)-dimensional
space.

Hamiltonian evolution traces a curve through phase space.

%%%%%%%%%%%%%%%%%%%%%%%%%%%%%%%%%%%%%%%%%%%%%%%%%%%%%%%%%%%%
\subsection{Poisson Brackets}
\label{subsec:poisson-brackets-physics}
%%%%%%%%%%%%%%%%%%%%%%%%%%%%%%%%%%%%%%%%%%%%%%%%%%%%%%%%%%%%

For functions \(f(q,p)\) and \(g(q,p)\), define the Poisson bracket

\[
\boxed{
\{f,g\}
=
\sum_i
\left(
\frac{\partial f}{\partial q_i}
\frac{\partial g}{\partial p_i}
-
\frac{\partial f}{\partial p_i}
\frac{\partial g}{\partial q_i}
\right).
}
\]

If \(f\) has no explicit time dependence,

\[
\boxed{
\frac{df}{dt}
=
\{f,H\}.
}
\]

Thus \(f\) is conserved when

\[
\{f,H\}=0.
\]

%%%%%%%%%%%%%%%%%%%%%%%%%%%%%%%%%%%%%%%%%%%%%%%%%%%%%%%%%%%%
\subsection{Symmetry}
\label{subsec:physical-symmetry}
%%%%%%%%%%%%%%%%%%%%%%%%%%%%%%%%%%%%%%%%%%%%%%%%%%%%%%%%%%%%

A symmetry is a transformation that leaves relevant physical structure
unchanged.

Examples include:

\[
\boxed{
\text{translation in space},
}
\]

\[
\boxed{
\text{translation in time},
}
\]

and

\[
\boxed{
\text{rotation}.
}
\]

Symmetry is deeply connected with conservation laws.

%%%%%%%%%%%%%%%%%%%%%%%%%%%%%%%%%%%%%%%%%%%%%%%%%%%%%%%%%%%%
\subsection{Noether's Theorem}
\label{subsec:noethers-theorem}
%%%%%%%%%%%%%%%%%%%%%%%%%%%%%%%%%%%%%%%%%%%%%%%%%%%%%%%%%%%%

Noether's theorem states, broadly, that continuous symmetries of the action
correspond to conserved quantities.

Important examples are:

\[
\boxed{
\text{time translation}
\longleftrightarrow
\text{energy conservation},
}
\]

\[
\boxed{
\text{spatial translation}
\longleftrightarrow
\text{linear momentum conservation},
}
\]

and

\[
\boxed{
\text{rotation}
\longleftrightarrow
\text{angular momentum conservation}.
}
\]

%%%%%%%%%%%%%%%%%%%%%%%%%%%%%%%%%%%%%%%%%%%%%%%%%%%%%%%%%%%%
\subsection{Scalar Fields}
\label{subsec:physical-scalar-fields}
%%%%%%%%%%%%%%%%%%%%%%%%%%%%%%%%%%%%%%%%%%%%%%%%%%%%%%%%%%%%

A scalar field assigns one number to each point in space and possibly time:

\[
\boxed{
\phi
=
\phi(\mathbf{x},t).
}
\]

Examples include:

\[
\boxed{
\text{temperature},
}
\]

\[
\boxed{
\text{pressure},
}
\]

and

\[
\boxed{
\text{electric potential}.
}
\]

%%%%%%%%%%%%%%%%%%%%%%%%%%%%%%%%%%%%%%%%%%%%%%%%%%%%%%%%%%%%
\subsection{Vector Fields}
\label{subsec:physical-vector-fields}
%%%%%%%%%%%%%%%%%%%%%%%%%%%%%%%%%%%%%%%%%%%%%%%%%%%%%%%%%%%%

A vector field assigns a vector to each point:

\[
\boxed{
\mathbf{F}
=
\mathbf{F}(\mathbf{x},t).
}
\]

Examples include:

\[
\boxed{
\text{velocity fields},
}
\]

\[
\boxed{
\text{electric fields},
}
\]

and

\[
\boxed{
\text{magnetic fields}.
}
\]

%%%%%%%%%%%%%%%%%%%%%%%%%%%%%%%%%%%%%%%%%%%%%%%%%%%%%%%%%%%%
\subsection{Gradient}
\label{subsec:physical-gradient}
%%%%%%%%%%%%%%%%%%%%%%%%%%%%%%%%%%%%%%%%%%%%%%%%%%%%%%%%%%%%

For scalar field \(\phi\),

\[
\boxed{
\nabla\phi
}
\]

points in the direction of greatest increase.

Its magnitude gives the maximum local rate of change.

Potential-generated forces often have form

\[
\boxed{
\mathbf{F}
=
-\nabla V.
}
\]

%%%%%%%%%%%%%%%%%%%%%%%%%%%%%%%%%%%%%%%%%%%%%%%%%%%%%%%%%%%%
\subsection{Divergence}
\label{subsec:physical-divergence}
%%%%%%%%%%%%%%%%%%%%%%%%%%%%%%%%%%%%%%%%%%%%%%%%%%%%%%%%%%%%

For vector field \(\mathbf{F}\),

\[
\boxed{
\nabla\cdot\mathbf{F}
}
\]

measures local source or sink behavior.

Positive divergence indicates net outward flux.

Negative divergence indicates net inward flux.

%%%%%%%%%%%%%%%%%%%%%%%%%%%%%%%%%%%%%%%%%%%%%%%%%%%%%%%%%%%%
\subsection{Curl}
\label{subsec:physical-curl}
%%%%%%%%%%%%%%%%%%%%%%%%%%%%%%%%%%%%%%%%%%%%%%%%%%%%%%%%%%%%

The curl

\[
\boxed{
\nabla\times\mathbf{F}
}
\]

measures local rotational structure of a vector field.

It appears naturally in fluid mechanics and electromagnetism.

%%%%%%%%%%%%%%%%%%%%%%%%%%%%%%%%%%%%%%%%%%%%%%%%%%%%%%%%%%%%
\subsection{Laplacian}
\label{subsec:physical-laplacian}
%%%%%%%%%%%%%%%%%%%%%%%%%%%%%%%%%%%%%%%%%%%%%%%%%%%%%%%%%%%%

For scalar field \(u\),

\[
\boxed{
\nabla^2u
=
\nabla\cdot\nabla u.
}
\]

In Cartesian coordinates,

\[
\boxed{
\nabla^2u
=
\frac{\partial^2u}{\partial x^2}
+
\frac{\partial^2u}{\partial y^2}
+
\frac{\partial^2u}{\partial z^2}.
}
\]

The Laplacian appears in diffusion, waves, electrostatics, and quantum
mechanics.

%%%%%%%%%%%%%%%%%%%%%%%%%%%%%%%%%%%%%%%%%%%%%%%%%%%%%%%%%%%%
\subsection{Conservation Law in Differential Form}
\label{subsec:conservation-law-differential-form}
%%%%%%%%%%%%%%%%%%%%%%%%%%%%%%%%%%%%%%%%%%%%%%%%%%%%%%%%%%%%

Let

\[
\rho(\mathbf{x},t)
\]

be a density and \(\mathbf{J}\) its flux.

Local conservation is written

\[
\boxed{
\frac{\partial\rho}{\partial t}
+
\nabla\cdot\mathbf{J}
=
0.
}
\]

This is the continuity equation.

%%%%%%%%%%%%%%%%%%%%%%%%%%%%%%%%%%%%%%%%%%%%%%%%%%%%%%%%%%%%
\subsection{Continuity Equation Interpretation}
\label{subsec:continuity-equation-interpretation}
%%%%%%%%%%%%%%%%%%%%%%%%%%%%%%%%%%%%%%%%%%%%%%%%%%%%%%%%%%%%

If

\[
\nabla\cdot\mathbf{J}>0,
\]

more quantity leaves a small region than enters.

Therefore density decreases.

The term

\[
-\nabla\cdot\mathbf{J}
\]

is the local rate of accumulation caused by transport.

%%%%%%%%%%%%%%%%%%%%%%%%%%%%%%%%%%%%%%%%%%%%%%%%%%%%%%%%%%%%
\subsection{Wave Equation}
\label{subsec:physical-wave-equation}
%%%%%%%%%%%%%%%%%%%%%%%%%%%%%%%%%%%%%%%%%%%%%%%%%%%%%%%%%%%%

The one-dimensional wave equation is

\[
\boxed{
\frac{\partial^2u}{\partial t^2}
=
c^2
\frac{\partial^2u}{\partial x^2},
}
\]

where \(c\) is wave speed.

It models idealized waves in strings, acoustics, electromagnetism, and other
systems.

%%%%%%%%%%%%%%%%%%%%%%%%%%%%%%%%%%%%%%%%%%%%%%%%%%%%%%%%%%%%
\subsection{Traveling Waves}
\label{subsec:traveling-waves}
%%%%%%%%%%%%%%%%%%%%%%%%%%%%%%%%%%%%%%%%%%%%%%%%%%%%%%%%%%%%

Functions of the form

\[
\boxed{
u(x,t)
=
f(x-ct)
}
\]

travel to the right with speed \(c\).

Similarly,

\[
\boxed{
u(x,t)
=
g(x+ct)
}
\]

travels to the left.

The general infinite-line solution can be written

\[
\boxed{
u(x,t)
=
f(x-ct)+g(x+ct).
}
\]

%%%%%%%%%%%%%%%%%%%%%%%%%%%%%%%%%%%%%%%%%%%%%%%%%%%%%%%%%%%%
\subsection{Worked Example: Verifying a Traveling Wave}
\label{subsec:worked-traveling-wave}
%%%%%%%%%%%%%%%%%%%%%%%%%%%%%%%%%%%%%%%%%%%%%%%%%%%%%%%%%%%%

\begin{workedexamplebox}{Wave Equation Verification}

Let

\[
u(x,t)
=
\sin(k(x-ct)).
\]

Then

\[
u_{tt}
=
-k^2c^2
\sin(k(x-ct)),
\]

and

\[
u_{xx}
=
-k^2
\sin(k(x-ct)).
\]

Hence

\[
u_{tt}
=
c^2u_{xx}.
\]

\begin{answerbox}
\[
\boxed{
u(x,t)
=
\sin(k(x-ct))
}
\]

solves the one-dimensional wave equation.
\end{answerbox}

\end{workedexamplebox}

%%%%%%%%%%%%%%%%%%%%%%%%%%%%%%%%%%%%%%%%%%%%%%%%%%%%%%%%%%%%
\subsection{Standing Waves}
\label{subsec:standing-waves}
%%%%%%%%%%%%%%%%%%%%%%%%%%%%%%%%%%%%%%%%%%%%%%%%%%%%%%%%%%%%

A standing wave may have form

\[
\boxed{
u(x,t)
=
X(x)T(t).
}
\]

For a string fixed at

\[
x=0,
\qquad
x=L,
\]

the spatial modes are

\[
\boxed{
X_n(x)
=
\sin
\left(
\frac{n\pi x}{L}
\right).
}
\]

The corresponding frequencies satisfy

\[
\boxed{
\omega_n
=
\frac{n\pi c}{L}.
}
\]

%%%%%%%%%%%%%%%%%%%%%%%%%%%%%%%%%%%%%%%%%%%%%%%%%%%%%%%%%%%%
\subsection{Normal Modes as Eigenfunctions}
\label{subsec:normal-modes-eigenfunctions}
%%%%%%%%%%%%%%%%%%%%%%%%%%%%%%%%%%%%%%%%%%%%%%%%%%%%%%%%%%%%

Separation of variables reduces the wave equation to an eigenvalue problem.

For example,

\[
\boxed{
-X''=\lambda X,
}
\]

with boundary conditions

\[
X(0)=X(L)=0.
\]

The eigenfunctions are

\[
\boxed{
X_n(x)
=
\sin
\left(
\frac{n\pi x}{L}
\right),
}
\]

with

\[
\boxed{
\lambda_n
=
\left(
\frac{n\pi}{L}
\right)^2.
}
\]

%%%%%%%%%%%%%%%%%%%%%%%%%%%%%%%%%%%%%%%%%%%%%%%%%%%%%%%%%%%%
\subsection{Heat Equation}
\label{subsec:physical-heat-equation}
%%%%%%%%%%%%%%%%%%%%%%%%%%%%%%%%%%%%%%%%%%%%%%%%%%%%%%%%%%%%

The one-dimensional heat equation is

\[
\boxed{
\frac{\partial u}{\partial t}
=
\alpha
\frac{\partial^2u}{\partial x^2},
}
\]

where \(\alpha>0\) is thermal diffusivity.

In several dimensions,

\[
\boxed{
u_t
=
\alpha\nabla^2u.
}
\]

The same mathematical structure describes diffusion.

%%%%%%%%%%%%%%%%%%%%%%%%%%%%%%%%%%%%%%%%%%%%%%%%%%%%%%%%%%%%
\subsection{Diffusion Smooths Gradients}
\label{subsec:diffusion-smoothing}
%%%%%%%%%%%%%%%%%%%%%%%%%%%%%%%%%%%%%%%%%%%%%%%%%%%%%%%%%%%%

If temperature is highly nonuniform, heat flows from hotter regions toward
cooler regions.

Mathematically, diffusion tends to reduce spatial variation.

This is reflected in the smoothing properties of the heat equation.

%%%%%%%%%%%%%%%%%%%%%%%%%%%%%%%%%%%%%%%%%%%%%%%%%%%%%%%%%%%%
\subsection{Laplace Equation}
\label{subsec:physical-laplace-equation}
%%%%%%%%%%%%%%%%%%%%%%%%%%%%%%%%%%%%%%%%%%%%%%%%%%%%%%%%%%%%

The Laplace equation is

\[
\boxed{
\nabla^2u=0.
}
\]

Solutions are called harmonic functions.

The equation appears in steady-state:

\[
\boxed{
\text{heat conduction},
}
\]

\[
\boxed{
\text{electrostatics},
}
\]

and

\[
\boxed{
\text{potential flow}.
}
\]

%%%%%%%%%%%%%%%%%%%%%%%%%%%%%%%%%%%%%%%%%%%%%%%%%%%%%%%%%%%%
\subsection{Poisson Equation}
\label{subsec:physical-poisson-equation}
%%%%%%%%%%%%%%%%%%%%%%%%%%%%%%%%%%%%%%%%%%%%%%%%%%%%%%%%%%%%

The Poisson equation is

\[
\boxed{
\nabla^2u=f.
}
\]

The source term \(f\) represents distributed forcing or density.

Examples include charge density in electrostatics and mass density in
Newtonian gravitational potential theory.

%%%%%%%%%%%%%%%%%%%%%%%%%%%%%%%%%%%%%%%%%%%%%%%%%%%%%%%%%%%%
\subsection{Boundary Conditions}
\label{subsec:physical-boundary-conditions}
%%%%%%%%%%%%%%%%%%%%%%%%%%%%%%%%%%%%%%%%%%%%%%%%%%%%%%%%%%%%

A PDE requires appropriate boundary conditions.

Common types include:

\[
\boxed{
u=g
}
\]

for Dirichlet conditions,

\[
\boxed{
\frac{\partial u}{\partial n}
=
h
}
\]

for Neumann conditions,

and

\[
\boxed{
au+b\frac{\partial u}{\partial n}=c
}
\]

for Robin conditions.

%%%%%%%%%%%%%%%%%%%%%%%%%%%%%%%%%%%%%%%%%%%%%%%%%%%%%%%%%%%%
\subsection{Initial Conditions}
\label{subsec:physical-initial-conditions}
%%%%%%%%%%%%%%%%%%%%%%%%%%%%%%%%%%%%%%%%%%%%%%%%%%%%%%%%%%%%

Time-dependent PDEs also require initial data.

For the heat equation,

\[
\boxed{
u(x,0)=f(x).
}
\]

For the wave equation, both position and velocity are typically required:

\[
\boxed{
u(x,0)=f(x),
}
\]

\[
\boxed{
u_t(x,0)=g(x).
}
\]

%%%%%%%%%%%%%%%%%%%%%%%%%%%%%%%%%%%%%%%%%%%%%%%%%%%%%%%%%%%%
\subsection{Fourier Series in Physics}
\label{subsec:fourier-series-physics}
%%%%%%%%%%%%%%%%%%%%%%%%%%%%%%%%%%%%%%%%%%%%%%%%%%%%%%%%%%%%

On an interval, a function may be expanded as

\[
\boxed{
f(x)
=
\sum_{n}
a_n\phi_n(x),
}
\]

where \(\phi_n\) are eigenfunctions adapted to the boundary conditions.

Each mode evolves independently in many linear PDEs.

This makes Fourier analysis one of the central mathematical tools of
physics.

%%%%%%%%%%%%%%%%%%%%%%%%%%%%%%%%%%%%%%%%%%%%%%%%%%%%%%%%%%%%
\subsection{Green's Functions}
\label{subsec:greens-functions-physics}
%%%%%%%%%%%%%%%%%%%%%%%%%%%%%%%%%%%%%%%%%%%%%%%%%%%%%%%%%%%%

For a linear differential operator \(L\), a Green's function \(G\) satisfies
schematically

\[
\boxed{
L G(x,\xi)
=
\delta(x-\xi),
}
\]

with appropriate boundary conditions.

Then solutions to

\[
L u=f
\]

may be written conceptually as

\[
\boxed{
u(x)
=
\int
G(x,\xi)f(\xi)\,d\xi.
}
\]

Green's functions represent the response to point forcing.

%%%%%%%%%%%%%%%%%%%%%%%%%%%%%%%%%%%%%%%%%%%%%%%%%%%%%%%%%%%%
\subsection{Electrostatics}
\label{subsec:electrostatics}
%%%%%%%%%%%%%%%%%%%%%%%%%%%%%%%%%%%%%%%%%%%%%%%%%%%%%%%%%%%%

The electric field is related to electric potential \(\phi\) by

\[
\boxed{
\mathbf{E}
=
-\nabla\phi.
}
\]

Gauss's law is

\[
\boxed{
\nabla\cdot\mathbf{E}
=
\frac{\rho}{\varepsilon_0}.
}
\]

Therefore,

\[
\boxed{
\nabla^2\phi
=
-\frac{\rho}{\varepsilon_0}.
}
\]

Thus electrostatic potential satisfies Poisson's equation.

%%%%%%%%%%%%%%%%%%%%%%%%%%%%%%%%%%%%%%%%%%%%%%%%%%%%%%%%%%%%
\subsection{Maxwell's Equations}
\label{subsec:maxwells-equations}
%%%%%%%%%%%%%%%%%%%%%%%%%%%%%%%%%%%%%%%%%%%%%%%%%%%%%%%%%%%%

In vacuum with charge density \(\rho\) and current density \(\mathbf{J}\),
Maxwell's equations can be written

\[
\boxed{
\nabla\cdot\mathbf{E}
=
\frac{\rho}{\varepsilon_0},
}
\]

\[
\boxed{
\nabla\cdot\mathbf{B}
=
0,
}
\]

\[
\boxed{
\nabla\times\mathbf{E}
=
-\frac{\partial\mathbf{B}}{\partial t},
}
\]

and

\[
\boxed{
\nabla\times\mathbf{B}
=
\mu_0\mathbf{J}
+
\mu_0\varepsilon_0
\frac{\partial\mathbf{E}}{\partial t}.
}
\]

These equations unify electricity, magnetism, and electromagnetic waves.

%%%%%%%%%%%%%%%%%%%%%%%%%%%%%%%%%%%%%%%%%%%%%%%%%%%%%%%%%%%%
\subsection{Electromagnetic Wave Speed}
\label{subsec:electromagnetic-wave-speed}
%%%%%%%%%%%%%%%%%%%%%%%%%%%%%%%%%%%%%%%%%%%%%%%%%%%%%%%%%%%%

In vacuum with

\[
\rho=0,
\qquad
\mathbf{J}=0,
\]

Maxwell's equations imply wave equations of the form

\[
\boxed{
\nabla^2\mathbf{E}
=
\mu_0\varepsilon_0
\frac{\partial^2\mathbf{E}}{\partial t^2}.
}
\]

Thus the wave speed is

\[
\boxed{
c
=
\frac{1}
{\sqrt{\mu_0\varepsilon_0}}.
}
\]

This identifies light as an electromagnetic wave.

%%%%%%%%%%%%%%%%%%%%%%%%%%%%%%%%%%%%%%%%%%%%%%%%%%%%%%%%%%%%
\subsection{Classical Field Theory}
\label{subsec:classical-field-theory}
%%%%%%%%%%%%%%%%%%%%%%%%%%%%%%%%%%%%%%%%%%%%%%%%%%%%%%%%%%%%

For fields

\[
\phi_a(\mathbf{x},t),
\]

one can define a Lagrangian density

\[
\boxed{
\mathcal{L}
=
\mathcal{L}
\left(
\phi_a,
\partial_\mu\phi_a,
x
\right).
}
\]

The action is

\[
\boxed{
S
=
\int
\mathcal{L}\,d^4x.
}
\]

Stationarity gives field Euler--Lagrange equations.

%%%%%%%%%%%%%%%%%%%%%%%%%%%%%%%%%%%%%%%%%%%%%%%%%%%%%%%%%%%%
\subsection{Field Euler--Lagrange Equations}
\label{subsec:field-euler-lagrange}
%%%%%%%%%%%%%%%%%%%%%%%%%%%%%%%%%%%%%%%%%%%%%%%%%%%%%%%%%%%%

For field \(\phi\),

\[
\boxed{
\frac{\partial\mathcal{L}}{\partial\phi}
-
\partial_\mu
\left(
\frac{\partial\mathcal{L}}
{\partial(\partial_\mu\phi)}
\right)
=
0.
}
\]

This generalizes particle mechanics to continuous fields.

%%%%%%%%%%%%%%%%%%%%%%%%%%%%%%%%%%%%%%%%%%%%%%%%%%%%%%%%%%%%
\subsection{Stress and Strain Preview}
\label{subsec:stress-strain-preview}
%%%%%%%%%%%%%%%%%%%%%%%%%%%%%%%%%%%%%%%%%%%%%%%%%%%%%%%%%%%%

Continuum mechanics treats matter as continuously distributed.

Stress describes internal forces per unit area.

Strain measures deformation.

A simple linear constitutive relation may be written schematically as

\[
\boxed{
\boldsymbol{\sigma}
=
\mathsf{C}:\boldsymbol{\varepsilon},
}
\]

where

\[
\boldsymbol{\sigma}
\]

is stress,

\[
\boldsymbol{\varepsilon}
\]

is strain, and

\[
\mathsf{C}
\]

is an elasticity tensor.

This is developed further in Section~13.4.

%%%%%%%%%%%%%%%%%%%%%%%%%%%%%%%%%%%%%%%%%%%%%%%%%%%%%%%%%%%%
\subsection{Special Relativity}
\label{subsec:special-relativity}
%%%%%%%%%%%%%%%%%%%%%%%%%%%%%%%%%%%%%%%%%%%%%%%%%%%%%%%%%%%%

Special relativity modifies classical space and time to preserve the speed
of light across inertial frames.

Define

\[
\boxed{
\gamma
=
\frac{1}
{\sqrt{1-v^2/c^2}}.
}
\]

The Lorentz transformation between frames moving at relative speed \(v\)
along the \(x\)-axis is

\[
\boxed{
x'
=
\gamma(x-vt),
}
\]

\[
\boxed{
t'
=
\gamma
\left(
t-\frac{vx}{c^2}
\right).
}
\]

%%%%%%%%%%%%%%%%%%%%%%%%%%%%%%%%%%%%%%%%%%%%%%%%%%%%%%%%%%%%
\subsection{Spacetime Interval}
\label{subsec:spacetime-interval}
%%%%%%%%%%%%%%%%%%%%%%%%%%%%%%%%%%%%%%%%%%%%%%%%%%%%%%%%%%%%

The quantity

\[
\boxed{
s^2
=
c^2t^2-x^2-y^2-z^2
}
\]

is invariant under Lorentz transformations under one common sign
convention.

This replaces ordinary Euclidean distance with a spacetime interval.

%%%%%%%%%%%%%%%%%%%%%%%%%%%%%%%%%%%%%%%%%%%%%%%%%%%%%%%%%%%%
\subsection{Time Dilation}
\label{subsec:time-dilation}
%%%%%%%%%%%%%%%%%%%%%%%%%%%%%%%%%%%%%%%%%%%%%%%%%%%%%%%%%%%%

If \(\Delta\tau\) is proper time measured in the rest frame of a clock, then
an inertial observer measuring the moving clock obtains

\[
\boxed{
\Delta t
=
\gamma\Delta\tau.
}
\]

Since

\[
\gamma\geq1,
\]

the moving clock is observed to accumulate less proper time between the same
events.

%%%%%%%%%%%%%%%%%%%%%%%%%%%%%%%%%%%%%%%%%%%%%%%%%%%%%%%%%%%%
\subsection{Relativistic Momentum}
\label{subsec:relativistic-momentum}
%%%%%%%%%%%%%%%%%%%%%%%%%%%%%%%%%%%%%%%%%%%%%%%%%%%%%%%%%%%%

Relativistic momentum is

\[
\boxed{
\mathbf{p}
=
\gamma m\mathbf{v}.
}
\]

It reduces to classical momentum

\[
m\mathbf{v}
\]

when

\[
v\ll c.
\]

%%%%%%%%%%%%%%%%%%%%%%%%%%%%%%%%%%%%%%%%%%%%%%%%%%%%%%%%%%%%
\subsection{Relativistic Energy}
\label{subsec:relativistic-energy}
%%%%%%%%%%%%%%%%%%%%%%%%%%%%%%%%%%%%%%%%%%%%%%%%%%%%%%%%%%%%

Total relativistic energy is

\[
\boxed{
E
=
\gamma mc^2.
}
\]

At rest,

\[
v=0,
\qquad
\gamma=1,
\]

giving

\[
\boxed{
E_0
=
mc^2.
}
\]

Energy and momentum satisfy

\[
\boxed{
E^2
=
p^2c^2
+
m^2c^4.
}
\]

%%%%%%%%%%%%%%%%%%%%%%%%%%%%%%%%%%%%%%%%%%%%%%%%%%%%%%%%%%%%
\subsection{General Relativity Preview}
\label{subsec:general-relativity-preview}
%%%%%%%%%%%%%%%%%%%%%%%%%%%%%%%%%%%%%%%%%%%%%%%%%%%%%%%%%%%%

General relativity models gravity geometrically.

Spacetime is described by a metric tensor

\[
\boxed{
g_{\mu\nu}.
}
\]

Matter and energy influence spacetime curvature, while free particles follow
geodesics determined by the metric.

The full theory uses differential geometry and tensor analysis.

%%%%%%%%%%%%%%%%%%%%%%%%%%%%%%%%%%%%%%%%%%%%%%%%%%%%%%%%%%%%
\subsection{Einstein Field Equation Preview}
\label{subsec:einstein-field-equation-preview}
%%%%%%%%%%%%%%%%%%%%%%%%%%%%%%%%%%%%%%%%%%%%%%%%%%%%%%%%%%%%

The Einstein field equations are written schematically as

\[
\boxed{
G_{\mu\nu}
+
\Lambda g_{\mu\nu}
=
\frac{8\pi G}{c^4}
T_{\mu\nu}.
}
\]

The left side describes spacetime geometry.

The tensor

\[
T_{\mu\nu}
\]

describes matter and energy.

This equation links geometry with physical content.

%%%%%%%%%%%%%%%%%%%%%%%%%%%%%%%%%%%%%%%%%%%%%%%%%%%%%%%%%%%%
\subsection{Quantum Mechanics}
\label{subsec:quantum-mechanics-mathematical-physics}
%%%%%%%%%%%%%%%%%%%%%%%%%%%%%%%%%%%%%%%%%%%%%%%%%%%%%%%%%%%%

Quantum mechanics represents a physical state by a vector or wavefunction in
a Hilbert space.

For a particle in one dimension,

\[
\boxed{
\psi(x,t)
}
\]

is the wavefunction.

The probability density for position is

\[
\boxed{
|\psi(x,t)|^2.
}
\]

%%%%%%%%%%%%%%%%%%%%%%%%%%%%%%%%%%%%%%%%%%%%%%%%%%%%%%%%%%%%
\subsection{Wavefunction Normalization}
\label{subsec:wavefunction-normalization}
%%%%%%%%%%%%%%%%%%%%%%%%%%%%%%%%%%%%%%%%%%%%%%%%%%%%%%%%%%%%

A physical wavefunction must satisfy

\[
\boxed{
\int_{-\infty}^{\infty}
|\psi(x,t)|^2\,dx
=
1.
}
\]

This ensures total probability equals \(1\).

%%%%%%%%%%%%%%%%%%%%%%%%%%%%%%%%%%%%%%%%%%%%%%%%%%%%%%%%%%%%
\subsection{Schrodinger Equation}
\label{subsec:schrodinger-equation}
%%%%%%%%%%%%%%%%%%%%%%%%%%%%%%%%%%%%%%%%%%%%%%%%%%%%%%%%%%%%

The time-dependent Schrödinger equation is

\[
\boxed{
i\hbar
\frac{\partial\psi}{\partial t}
=
\hat H\psi.
}
\]

For a nonrelativistic particle in one dimension,

\[
\boxed{
\hat H
=
-\frac{\hbar^2}{2m}
\frac{\partial^2}{\partial x^2}
+
V(x).
}
\]

Therefore,

\[
\boxed{
i\hbar\psi_t
=
-\frac{\hbar^2}{2m}\psi_{xx}
+
V\psi.
}
\]

%%%%%%%%%%%%%%%%%%%%%%%%%%%%%%%%%%%%%%%%%%%%%%%%%%%%%%%%%%%%
\subsection{Stationary States}
\label{subsec:quantum-stationary-states}
%%%%%%%%%%%%%%%%%%%%%%%%%%%%%%%%%%%%%%%%%%%%%%%%%%%%%%%%%%%%

If \(V\) is independent of time, seek

\[
\boxed{
\psi(x,t)
=
\phi(x)e^{-iEt/\hbar}.
}
\]

Substitution gives the time-independent Schrödinger equation

\[
\boxed{
\hat H\phi
=
E\phi.
}
\]

Thus stationary-state energies are eigenvalues of the Hamiltonian operator.

%%%%%%%%%%%%%%%%%%%%%%%%%%%%%%%%%%%%%%%%%%%%%%%%%%%%%%%%%%%%
\subsection{Quantum Observables}
\label{subsec:quantum-observables}
%%%%%%%%%%%%%%%%%%%%%%%%%%%%%%%%%%%%%%%%%%%%%%%%%%%%%%%%%%%%

Physical observables are represented by suitable linear operators.

For example, position is

\[
\boxed{
\hat{x}\psi
=
x\psi,
}
\]

while momentum in one dimension is

\[
\boxed{
\hat{p}
=
-i\hbar\frac{d}{dx}.
}
\]

Possible measurement values are associated with operator spectra.

%%%%%%%%%%%%%%%%%%%%%%%%%%%%%%%%%%%%%%%%%%%%%%%%%%%%%%%%%%%%
\subsection{Expectation Values}
\label{subsec:quantum-expectation-values}
%%%%%%%%%%%%%%%%%%%%%%%%%%%%%%%%%%%%%%%%%%%%%%%%%%%%%%%%%%%%

For observable \(\hat A\), the expectation value is

\[
\boxed{
\langle A\rangle
=
\int
\psi^*
\hat A
\psi
\,dx,
}
\]

assuming appropriate normalization and domain conditions.

This gives the statistical mean over repeated measurements of identically
prepared systems.

%%%%%%%%%%%%%%%%%%%%%%%%%%%%%%%%%%%%%%%%%%%%%%%%%%%%%%%%%%%%
\subsection{Commutators}
\label{subsec:quantum-commutators}
%%%%%%%%%%%%%%%%%%%%%%%%%%%%%%%%%%%%%%%%%%%%%%%%%%%%%%%%%%%%

For operators \(\hat A\) and \(\hat B\), define

\[
\boxed{
[\hat A,\hat B]
=
\hat A\hat B-\hat B\hat A.
}
\]

Position and momentum satisfy

\[
\boxed{
[\hat x,\hat p]
=
i\hbar.
}
\]

Noncommutativity is fundamental to quantum mechanics.

%%%%%%%%%%%%%%%%%%%%%%%%%%%%%%%%%%%%%%%%%%%%%%%%%%%%%%%%%%%%
\subsection{Uncertainty Relation}
\label{subsec:quantum-uncertainty}
%%%%%%%%%%%%%%%%%%%%%%%%%%%%%%%%%%%%%%%%%%%%%%%%%%%%%%%%%%%%

Position and momentum uncertainties satisfy

\[
\boxed{
\Delta x\,\Delta p
\geq
\frac{\hbar}{2}.
}
\]

This is not merely experimental imperfection.

It follows from the mathematical structure of noncommuting observables.

%%%%%%%%%%%%%%%%%%%%%%%%%%%%%%%%%%%%%%%%%%%%%%%%%%%%%%%%%%%%
\subsection{Infinite Square Well}
\label{subsec:infinite-square-well}
%%%%%%%%%%%%%%%%%%%%%%%%%%%%%%%%%%%%%%%%%%%%%%%%%%%%%%%%%%%%

Consider

\[
V(x)=0
\]

for

\[
0<x<L,
\]

with infinite barriers outside.

Boundary conditions are

\[
\boxed{
\phi(0)=0,
\qquad
\phi(L)=0.
}
\]

Inside the well,

\[
-\frac{\hbar^2}{2m}\phi''
=
E\phi.
\]

%%%%%%%%%%%%%%%%%%%%%%%%%%%%%%%%%%%%%%%%%%%%%%%%%%%%%%%%%%%%
\subsection{Energy Levels of the Infinite Well}
\label{subsec:infinite-well-energy-levels}
%%%%%%%%%%%%%%%%%%%%%%%%%%%%%%%%%%%%%%%%%%%%%%%%%%%%%%%%%%%%

The normalized eigenfunctions are

\[
\boxed{
\phi_n(x)
=
\sqrt{\frac{2}{L}}
\sin
\left(
\frac{n\pi x}{L}
\right),
}
\]

for

\[
n=1,2,3,\ldots.
\]

The energy levels are

\[
\boxed{
E_n
=
\frac{n^2\pi^2\hbar^2}
{2mL^2}.
}
\]

The energy is quantized.

%%%%%%%%%%%%%%%%%%%%%%%%%%%%%%%%%%%%%%%%%%%%%%%%%%%%%%%%%%%%
\subsection{Worked Example: Ground-State Energy}
\label{subsec:worked-ground-state-energy}
%%%%%%%%%%%%%%%%%%%%%%%%%%%%%%%%%%%%%%%%%%%%%%%%%%%%%%%%%%%%

\begin{workedexamplebox}{Lowest Energy in a Box}

For the infinite square well,

\[
E_n
=
\frac{n^2\pi^2\hbar^2}
{2mL^2}.
\]

The ground state corresponds to

\[
n=1.
\]

Therefore,

\begin{answerbox}
\[
\boxed{
E_1
=
\frac{\pi^2\hbar^2}
{2mL^2}.
}
\]
\end{answerbox}

The lowest possible energy is not zero.

\end{workedexamplebox}

%%%%%%%%%%%%%%%%%%%%%%%%%%%%%%%%%%%%%%%%%%%%%%%%%%%%%%%%%%%%
\subsection{Quantum Superposition}
\label{subsec:quantum-superposition}
%%%%%%%%%%%%%%%%%%%%%%%%%%%%%%%%%%%%%%%%%%%%%%%%%%%%%%%%%%%%

If

\[
\psi_1
\]

and

\[
\psi_2
\]

are valid states, then a linear combination

\[
\boxed{
\psi
=
c_1\psi_1+c_2\psi_2
}
\]

is also a valid state when properly normalized.

This follows from linearity of the Schrödinger equation.

%%%%%%%%%%%%%%%%%%%%%%%%%%%%%%%%%%%%%%%%%%%%%%%%%%%%%%%%%%%%
\subsection{Hilbert Space}
\label{subsec:quantum-hilbert-space}
%%%%%%%%%%%%%%%%%%%%%%%%%%%%%%%%%%%%%%%%%%%%%%%%%%%%%%%%%%%%

Quantum states belong to a Hilbert space.

This provides:

\[
\boxed{
\text{vector addition},
}
\]

\[
\boxed{
\text{scalar multiplication},
}
\]

\[
\boxed{
\text{inner products},
}
\]

and

\[
\boxed{
\text{completeness}.
}
\]

Linear algebra and functional analysis therefore form the mathematical
language of quantum theory.

%%%%%%%%%%%%%%%%%%%%%%%%%%%%%%%%%%%%%%%%%%%%%%%%%%%%%%%%%%%%
\subsection{Quantum Harmonic Oscillator}
\label{subsec:quantum-harmonic-oscillator}
%%%%%%%%%%%%%%%%%%%%%%%%%%%%%%%%%%%%%%%%%%%%%%%%%%%%%%%%%%%%

The potential is

\[
\boxed{
V(x)
=
\frac12m\omega^2x^2.
}
\]

The energy eigenvalues are

\[
\boxed{
E_n
=
\hbar\omega
\left(
n+\frac12
\right),
}
\]

for

\[
n=0,1,2,\ldots.
\]

The minimum energy is

\[
\boxed{
E_0
=
\frac12\hbar\omega.
}
\]

%%%%%%%%%%%%%%%%%%%%%%%%%%%%%%%%%%%%%%%%%%%%%%%%%%%%%%%%%%%%
\subsection{Classical and Quantum Comparison}
\label{subsec:classical-quantum-comparison}
%%%%%%%%%%%%%%%%%%%%%%%%%%%%%%%%%%%%%%%%%%%%%%%%%%%%%%%%%%%%

Classical mechanics describes a particle state through definite coordinates
and momenta.

Quantum mechanics describes a state by a Hilbert-space vector.

Classical observables are functions such as

\[
f(q,p).
\]

Quantum observables become operators.

This correspondence motivates the transition from Poisson brackets to
operator commutators.

%%%%%%%%%%%%%%%%%%%%%%%%%%%%%%%%%%%%%%%%%%%%%%%%%%%%%%%%%%%%
\subsection{Dimensional Analysis in Physics}
\label{subsec:dimensional-analysis-physics}
%%%%%%%%%%%%%%%%%%%%%%%%%%%%%%%%%%%%%%%%%%%%%%%%%%%%%%%%%%%%

Suppose the period \(T\) of a simple pendulum depends only on length \(L\)
and gravitational acceleration \(g\).

Assume

\[
T
=
C L^a g^b.
\]

Dimensions give

\[
[T]
=
L^a
(LT^{-2})^b.
\]

Thus

\[
T
=
L^{a+b}
T^{-2b}.
\]

Matching exponents gives

\[
a+b=0,
\]

\[
-2b=1.
\]

Hence

\[
b=-\frac12,
\qquad
a=\frac12.
\]

Therefore,

\[
\boxed{
T
\propto
\sqrt{\frac{L}{g}}.
}
\]

%%%%%%%%%%%%%%%%%%%%%%%%%%%%%%%%%%%%%%%%%%%%%%%%%%%%%%%%%%%%
\subsection{Worked Example: Pendulum Scaling}
\label{subsec:worked-pendulum-scaling}
%%%%%%%%%%%%%%%%%%%%%%%%%%%%%%%%%%%%%%%%%%%%%%%%%%%%%%%%%%%%

\begin{workedexamplebox}{Finding the Functional Dependence}

From dimensional analysis,

\[
T
=
C
\sqrt{\frac{L}{g}}.
\]

Dimensional reasoning cannot determine dimensionless constant \(C\).

For the small-angle pendulum equation, solving the dynamics gives

\[
C=2\pi.
\]

Thus

\begin{answerbox}
\[
\boxed{
T
=
2\pi
\sqrt{\frac{L}{g}}.
}
\]
\end{answerbox}

\end{workedexamplebox}

%%%%%%%%%%%%%%%%%%%%%%%%%%%%%%%%%%%%%%%%%%%%%%%%%%%%%%%%%%%%
\subsection{Nondimensionalization in Physics}
\label{subsec:nondimensionalization-physics}
%%%%%%%%%%%%%%%%%%%%%%%%%%%%%%%%%%%%%%%%%%%%%%%%%%%%%%%%%%%%

Consider

\[
m\ddot{x}
+
c\dot{x}
+
kx
=
0.
\]

Choose characteristic time

\[
\boxed{
T
=
\sqrt{\frac{m}{k}}.
}
\]

Let

\[
\tau
=
\frac{t}{T}.
\]

Then the equation becomes

\[
\boxed{
\frac{d^2x}{d\tau^2}
+
2\zeta
\frac{dx}{d\tau}
+
x
=
0,
}
\]

where

\[
\boxed{
\zeta
=
\frac{c}{2\sqrt{mk}}
}
\]

is the dimensionless damping ratio.

%%%%%%%%%%%%%%%%%%%%%%%%%%%%%%%%%%%%%%%%%%%%%%%%%%%%%%%%%%%%
\subsection{Dimensionless Numbers}
\label{subsec:physical-dimensionless-numbers}
%%%%%%%%%%%%%%%%%%%%%%%%%%%%%%%%%%%%%%%%%%%%%%%%%%%%%%%%%%%%

Dimensionless combinations often determine regimes of physical behavior.

Examples include:

\[
\boxed{
\text{Reynolds number},
}
\]

\[
\boxed{
\text{Mach number},
}
\]

\[
\boxed{
\text{Peclet number},
}
\]

and

\[
\boxed{
\text{Froude number}.
}
\]

Such quantities compare competing physical effects.

%%%%%%%%%%%%%%%%%%%%%%%%%%%%%%%%%%%%%%%%%%%%%%%%%%%%%%%%%%%%
\subsection{Perturbation Methods}
\label{subsec:physical-perturbation-methods}
%%%%%%%%%%%%%%%%%%%%%%%%%%%%%%%%%%%%%%%%%%%%%%%%%%%%%%%%%%%%

Suppose a problem depends on a small parameter

\[
0<\varepsilon\ll1.
\]

One may seek an expansion

\[
\boxed{
x
=
x_0
+
\varepsilon x_1
+
\varepsilon^2x_2
+
\cdots.
}
\]

Substitution into the governing equation gives a hierarchy of simpler
problems.

%%%%%%%%%%%%%%%%%%%%%%%%%%%%%%%%%%%%%%%%%%%%%%%%%%%%%%%%%%%%
\subsection{Regular Perturbation Expansion}
\label{subsec:regular-perturbation}
%%%%%%%%%%%%%%%%%%%%%%%%%%%%%%%%%%%%%%%%%%%%%%%%%%%%%%%%%%%%

Consider

\[
F(x,\varepsilon)=0.
\]

Assume

\[
x=x_0+\varepsilon x_1+\cdots.
\]

Matching powers of \(\varepsilon\) produces equations for

\[
x_0,
\quad
x_1,
\quad
\ldots.
\]

The approximation is valid only when the perturbation expansion remains
uniformly meaningful in the region of interest.

%%%%%%%%%%%%%%%%%%%%%%%%%%%%%%%%%%%%%%%%%%%%%%%%%%%%%%%%%%%%
\subsection{Singular Perturbations}
\label{subsec:singular-perturbations-physics}
%%%%%%%%%%%%%%%%%%%%%%%%%%%%%%%%%%%%%%%%%%%%%%%%%%%%%%%%%%%%

A perturbation is singular when setting

\[
\varepsilon=0
\]

changes the character or order of the governing problem.

For example,

\[
\boxed{
\varepsilon y''
+
y'
+
y
=
0
}
\]

loses one derivative when \(\varepsilon=0\).

Boundary layers and multiple scales commonly arise in such systems.

%%%%%%%%%%%%%%%%%%%%%%%%%%%%%%%%%%%%%%%%%%%%%%%%%%%%%%%%%%%%
\subsection{Asymptotic Approximation}
\label{subsec:physical-asymptotics}
%%%%%%%%%%%%%%%%%%%%%%%%%%%%%%%%%%%%%%%%%%%%%%%%%%%%%%%%%%%%

An asymptotic approximation describes behavior in a limiting regime.

For example,

\[
f(x)
\sim
g(x)
\]

as

\[
x\to\infty
\]

means

\[
\boxed{
\frac{f(x)}{g(x)}
\to1.
}
\]

Asymptotic methods are central when exact solutions are unavailable.

%%%%%%%%%%%%%%%%%%%%%%%%%%%%%%%%%%%%%%%%%%%%%%%%%%%%%%%%%%%%
\subsection{Variational Methods}
\label{subsec:variational-methods-physics}
%%%%%%%%%%%%%%%%%%%%%%%%%%%%%%%%%%%%%%%%%%%%%%%%%%%%%%%%%%%%

Many physical problems can be formulated as minimizing or stationarizing a
functional.

For example,

\[
\boxed{
J[u]
=
\int_\Omega
F
\left(
x,u,\nabla u
\right)
dx.
}
\]

Stationarity conditions lead to differential equations.

Variational formulations are especially useful in mechanics and PDE theory.

%%%%%%%%%%%%%%%%%%%%%%%%%%%%%%%%%%%%%%%%%%%%%%%%%%%%%%%%%%%%
\subsection{Energy Minimization}
\label{subsec:physical-energy-minimization}
%%%%%%%%%%%%%%%%%%%%%%%%%%%%%%%%%%%%%%%%%%%%%%%%%%%%%%%%%%%%

Stable equilibrium configurations often minimize an energy functional.

For example, an elastic system may seek

\[
\boxed{
u^*
=
\arg\min_u
E[u].
}
\]

This connects physical equilibrium with optimization and calculus of
variations.

%%%%%%%%%%%%%%%%%%%%%%%%%%%%%%%%%%%%%%%%%%%%%%%%%%%%%%%%%%%%
\subsection{Weak Formulations}
\label{subsec:weak-formulations-physics}
%%%%%%%%%%%%%%%%%%%%%%%%%%%%%%%%%%%%%%%%%%%%%%%%%%%%%%%%%%%%

A differential equation can sometimes be rewritten by multiplying by a test
function and integrating.

Derivatives may then be transferred using integration by parts.

This produces a weak formulation requiring fewer classical derivatives.

Weak formulations are foundational for:

\[
\boxed{
\text{variational PDE theory}
}
\]

and

\[
\boxed{
\text{finite element methods}.
}
\]

%%%%%%%%%%%%%%%%%%%%%%%%%%%%%%%%%%%%%%%%%%%%%%%%%%%%%%%%%%%%
\subsection{Numerical Mathematical Physics}
\label{subsec:numerical-mathematical-physics}
%%%%%%%%%%%%%%%%%%%%%%%%%%%%%%%%%%%%%%%%%%%%%%%%%%%%%%%%%%%%

Many physically important equations have no closed-form solutions.

Numerical methods approximate:

\[
\boxed{
\text{particle trajectories},
}
\]

\[
\boxed{
\text{wave fields},
}
\]

\[
\boxed{
\text{temperature fields},
}
\]

\[
\boxed{
\text{fluid flows},
}
\]

and

\[
\boxed{
\text{quantum states}.
}
\]

Simulation is therefore central to modern mathematical physics.

%%%%%%%%%%%%%%%%%%%%%%%%%%%%%%%%%%%%%%%%%%%%%%%%%%%%%%%%%%%%
\subsection{Time-Stepping}
\label{subsec:physical-time-stepping}
%%%%%%%%%%%%%%%%%%%%%%%%%%%%%%%%%%%%%%%%%%%%%%%%%%%%%%%%%%%%

For

\[
\dot{x}=f(x,t),
\]

a simple explicit Euler step is

\[
\boxed{
x_{n+1}
=
x_n
+
\Delta t
f(x_n,t_n).
}
\]

More sophisticated methods improve accuracy and stability.

The choice of numerical method should reflect the physical and mathematical
structure of the problem.

%%%%%%%%%%%%%%%%%%%%%%%%%%%%%%%%%%%%%%%%%%%%%%%%%%%%%%%%%%%%
\subsection{Conservation and Numerical Methods}
\label{subsec:numerical-conservation}
%%%%%%%%%%%%%%%%%%%%%%%%%%%%%%%%%%%%%%%%%%%%%%%%%%%%%%%%%%%%

A numerical scheme may introduce artificial drift in quantities that the
continuous system conserves.

For example, a poorly chosen integrator can cause energy to grow or decay
artificially in a conservative oscillator.

Structure-preserving methods attempt to respect properties such as:

\[
\boxed{
\text{energy},
}
\]

\[
\boxed{
\text{momentum},
}
\]

or

\[
\boxed{
\text{symplectic geometry}.
}
\]

%%%%%%%%%%%%%%%%%%%%%%%%%%%%%%%%%%%%%%%%%%%%%%%%%%%%%%%%%%%%
\subsection{Mathematical Error Versus Numerical Error}
\label{subsec:physical-model-numerical-error}
%%%%%%%%%%%%%%%%%%%%%%%%%%%%%%%%%%%%%%%%%%%%%%%%%%%%%%%%%%%%

Two errors should be distinguished.

\[
\boxed{
\text{model error}
}
\]

arises because the mathematical equations approximate physical reality.

\[
\boxed{
\text{numerical error}
}
\]

arises because the equations themselves are solved approximately.

Improving the numerical solver does not remove model error.

%%%%%%%%%%%%%%%%%%%%%%%%%%%%%%%%%%%%%%%%%%%%%%%%%%%%%%%%%%%%
\subsection{Mathematical Physics Workflow}
\label{subsec:mathematical-physics-workflow}
%%%%%%%%%%%%%%%%%%%%%%%%%%%%%%%%%%%%%%%%%%%%%%%%%%%%%%%%%%%%

A useful workflow is:

\begin{enumerate}
    \item define the physical system;
    \item identify relevant scales;
    \item choose state variables and fields;
    \item list physical parameters and units;
    \item identify symmetries;
    \item identify conserved quantities;
    \item determine whether a particle or continuum model is appropriate;
    \item choose governing physical laws;
    \item translate those laws into equations;
    \item verify dimensional consistency;
    \item determine initial and boundary conditions;
    \item nondimensionalize when useful;
    \item identify important dimensionless parameters;
    \item search for equilibria or steady states;
    \item analyze stability;
    \item identify normal modes or eigenvalue problems;
    \item exploit conservation laws;
    \item exploit symmetry;
    \item use variational principles when appropriate;
    \item determine whether exact solutions exist;
    \item use perturbation or asymptotic methods when a small parameter
          exists;
    \item use numerical methods when analytic methods are insufficient;
    \item check convergence of the numerical approximation;
    \item compare model predictions with physical observations;
    \item distinguish model error from computational error.
\end{enumerate}

%%%%%%%%%%%%%%%%%%%%%%%%%%%%%%%%%%%%%%%%%%%%%%%%%%%%%%%%%%%%
\subsection{Common Errors}
\label{subsec:mathematical-physics-common-errors}
%%%%%%%%%%%%%%%%%%%%%%%%%%%%%%%%%%%%%%%%%%%%%%%%%%%%%%%%%%%%

\subsubsection{Ignoring Units}

Equations that add quantities of different dimensions are physically
invalid.

Dimensional consistency should be checked before detailed computation.

%%%%%%%%%%%%%%%%%%%%%%%%%%%%%%%%%%%%%%%%%%%%%%%%%%%%%%%%%%%%
\subsubsection{Confusing Position, Velocity, and Acceleration}

These quantities satisfy

\[
\mathbf{v}
=
\dot{\mathbf{r}},
\]

\[
\mathbf{a}
=
\dot{\mathbf{v}}.
\]

They have different dimensions and physical interpretations.

%%%%%%%%%%%%%%%%%%%%%%%%%%%%%%%%%%%%%%%%%%%%%%%%%%%%%%%%%%%%
\subsubsection{Using \(F=ma\) Without Checking the Model}

The form

\[
F=ma
\]

assumes constant mass and an inertial reference frame in ordinary Newtonian
mechanics.

More general systems may require momentum balance or relativistic dynamics.

%%%%%%%%%%%%%%%%%%%%%%%%%%%%%%%%%%%%%%%%%%%%%%%%%%%%%%%%%%%%
\subsubsection{Confusing Potential With Force}

Potential energy is scalar.

Force is

\[
\boxed{
\mathbf{F}
=
-\nabla V.
}
\]

They are not the same quantity.

%%%%%%%%%%%%%%%%%%%%%%%%%%%%%%%%%%%%%%%%%%%%%%%%%%%%%%%%%%%%
\subsubsection{Forgetting the Sign in Conservative Forces}

The negative sign matters:

\[
\boxed{
\mathbf{F}
=
-\nabla V.
}
\]

The force points toward decreasing potential energy.

%%%%%%%%%%%%%%%%%%%%%%%%%%%%%%%%%%%%%%%%%%%%%%%%%%%%%%%%%%%%
\subsubsection{Assuming Energy Is Always Conserved}

Mechanical energy is conserved under appropriate conservative,
time-independent conditions.

Damping, external forcing, or explicit time dependence can change energy.

%%%%%%%%%%%%%%%%%%%%%%%%%%%%%%%%%%%%%%%%%%%%%%%%%%%%%%%%%%%%
\subsubsection{Confusing Angular Momentum and Torque}

Angular momentum is

\[
\mathbf{L}
=
\mathbf{r}\times\mathbf{p}.
\]

Torque is its time derivative:

\[
\boldsymbol{\tau}
=
\frac{d\mathbf{L}}{dt}.
\]

%%%%%%%%%%%%%%%%%%%%%%%%%%%%%%%%%%%%%%%%%%%%%%%%%%%%%%%%%%%%
\subsubsection{Assuming Every Oscillation Is Simple Harmonic Motion}

Simple harmonic motion requires a restoring force linear in displacement.

Nonlinear oscillators may have amplitude-dependent frequencies and more
complicated dynamics.

%%%%%%%%%%%%%%%%%%%%%%%%%%%%%%%%%%%%%%%%%%%%%%%%%%%%%%%%%%%%
\subsubsection{Confusing Natural and Driving Frequencies}

The natural frequency is a property of the unforced system.

The driving frequency comes from external forcing.

Resonance depends on their relationship.

%%%%%%%%%%%%%%%%%%%%%%%%%%%%%%%%%%%%%%%%%%%%%%%%%%%%%%%%%%%%
\subsubsection{Assuming Resonance Requires Zero Damping}

Damped systems can still exhibit resonant amplification.

Damping changes the resonance peak and prevents the ideal unbounded growth
of the undamped resonantly forced model.

%%%%%%%%%%%%%%%%%%%%%%%%%%%%%%%%%%%%%%%%%%%%%%%%%%%%%%%%%%%%
\subsubsection{Confusing Generalized Coordinates With Cartesian Coordinates}

Generalized coordinates may be:

\[
\boxed{
\text{angles},
}
\]

\[
\boxed{
\text{distances},
}
\]

or other parameters describing the configuration.

They need not be ordinary Cartesian positions.

%%%%%%%%%%%%%%%%%%%%%%%%%%%%%%%%%%%%%%%%%%%%%%%%%%%%%%%%%%%%
\subsubsection{Assuming the Hamiltonian Is Always Equal to Total Energy}

For many standard conservative mechanical systems,

\[
H=T+V.
\]

But in more general coordinate systems or explicitly time-dependent systems,
the interpretation requires care.

%%%%%%%%%%%%%%%%%%%%%%%%%%%%%%%%%%%%%%%%%%%%%%%%%%%%%%%%%%%%
\subsubsection{Confusing Phase Space With Physical Space}

Physical space contains positions.

Phase space contains both coordinates and momenta.

For \(n\) degrees of freedom, phase space typically has dimension \(2n\).

%%%%%%%%%%%%%%%%%%%%%%%%%%%%%%%%%%%%%%%%%%%%%%%%%%%%%%%%%%%%
\subsubsection{Confusing Divergence and Curl}

Divergence measures local source behavior.

Curl measures local rotational behavior.

They encode different geometric information.

%%%%%%%%%%%%%%%%%%%%%%%%%%%%%%%%%%%%%%%%%%%%%%%%%%%%%%%%%%%%
\subsubsection{Ignoring Boundary Conditions}

A PDE alone usually does not determine a unique physical solution.

Boundary and initial conditions are essential parts of the model.

%%%%%%%%%%%%%%%%%%%%%%%%%%%%%%%%%%%%%%%%%%%%%%%%%%%%%%%%%%%%
\subsubsection{Confusing the Heat and Wave Equations}

The heat equation is first order in time:

\[
u_t=\alpha\nabla^2u.
\]

The wave equation is second order in time:

\[
u_{tt}=c^2\nabla^2u.
\]

Their qualitative behavior is very different.

%%%%%%%%%%%%%%%%%%%%%%%%%%%%%%%%%%%%%%%%%%%%%%%%%%%%%%%%%%%%
\subsubsection{Assuming All PDE Solutions Are Smooth}

Physical and mathematical solutions may develop discontinuities,
singularities, or insufficient regularity for classical derivatives.

Weak or generalized solutions may be required.

%%%%%%%%%%%%%%%%%%%%%%%%%%%%%%%%%%%%%%%%%%%%%%%%%%%%%%%%%%%%
\subsubsection{Confusing Electric Potential and Electric Field}

They are related by

\[
\boxed{
\mathbf{E}
=
-\nabla\phi.
}
\]

The potential is scalar.

The electric field is vector-valued.

%%%%%%%%%%%%%%%%%%%%%%%%%%%%%%%%%%%%%%%%%%%%%%%%%%%%%%%%%%%%
\subsubsection{Treating Relativistic Formulas as Classical Corrections at All Speeds}

Relativistic mechanics replaces, rather than merely slightly adjusts,
classical relationships at speeds comparable to \(c\).

Classical formulas arise as low-speed approximations.

%%%%%%%%%%%%%%%%%%%%%%%%%%%%%%%%%%%%%%%%%%%%%%%%%%%%%%%%%%%%
\subsubsection{Confusing Proper Time and Coordinate Time}

Proper time is measured along the worldline of a clock.

Coordinate time depends on the chosen reference frame.

They are related but not identical.

%%%%%%%%%%%%%%%%%%%%%%%%%%%%%%%%%%%%%%%%%%%%%%%%%%%%%%%%%%%%
\subsubsection{Interpreting the Wavefunction as an Ordinary Probability Density}

The wavefunction \(\psi\) can be complex.

The probability density is

\[
\boxed{
|\psi|^2,
}
\]

not \(\psi\) itself.

%%%%%%%%%%%%%%%%%%%%%%%%%%%%%%%%%%%%%%%%%%%%%%%%%%%%%%%%%%%%
\subsubsection{Forgetting Quantum Normalization}

A physical state should be normalized so that

\[
\int |\psi|^2\,dx=1.
\]

Otherwise probabilities are not properly normalized.

%%%%%%%%%%%%%%%%%%%%%%%%%%%%%%%%%%%%%%%%%%%%%%%%%%%%%%%%%%%%
\subsubsection{Confusing an Operator With Its Eigenvalue}

An observable is represented by an operator.

Its eigenvalues are possible measurement values in appropriate states.

%%%%%%%%%%%%%%%%%%%%%%%%%%%%%%%%%%%%%%%%%%%%%%%%%%%%%%%%%%%%
\subsubsection{Assuming Quantum Uncertainty Is Only Experimental Noise}

Quantum uncertainty arises mathematically from the structure of the state
and noncommuting observables.

It is not simply poor measurement equipment.

%%%%%%%%%%%%%%%%%%%%%%%%%%%%%%%%%%%%%%%%%%%%%%%%%%%%%%%%%%%%
\subsubsection{Confusing Quantization With Numerical Discretization}

Quantum energy levels can be mathematically discrete because of spectral and
boundary conditions.

This is different from approximating a continuous problem using a numerical
grid.

%%%%%%%%%%%%%%%%%%%%%%%%%%%%%%%%%%%%%%%%%%%%%%%%%%%%%%%%%%%%
\subsubsection{Ignoring Small Parameters in Physical Models}

A small dimensionless parameter can indicate a useful approximation or
separation of scales.

Nondimensionalization often reveals these parameters.

%%%%%%%%%%%%%%%%%%%%%%%%%%%%%%%%%%%%%%%%%%%%%%%%%%%%%%%%%%%%
\subsubsection{Assuming a Perturbation Expansion Is Always Valid}

A formal series may fail near boundaries, resonances, long time scales, or
singular limits.

Its domain of validity must be analyzed.

%%%%%%%%%%%%%%%%%%%%%%%%%%%%%%%%%%%%%%%%%%%%%%%%%%%%%%%%%%%%
\subsubsection{Confusing Numerical Accuracy With Physical Accuracy}

A numerical solution may approximate the chosen equations extremely well
while the equations themselves poorly represent the physical system.

%%%%%%%%%%%%%%%%%%%%%%%%%%%%%%%%%%%%%%%%%%%%%%%%%%%%%%%%%%%%
\subsection{Applications}
\label{subsec:mathematical-physics-applications}
%%%%%%%%%%%%%%%%%%%%%%%%%%%%%%%%%%%%%%%%%%%%%%%%%%%%%%%%%%%%

\begin{applicationbox}

\textbf{Classical Mechanics.}
Differential equations, conservation laws, Lagrangian mechanics, and
Hamiltonian systems describe particle and rigid-body motion.

\medskip

\textbf{Oscillations and Waves.}
Eigenvalue methods, Fourier analysis, and PDEs describe vibrating
structures, acoustics, optics, and wave propagation.

\medskip

\textbf{Electromagnetism.}
Vector calculus and PDEs organize electric and magnetic fields through
Maxwell's equations.

\medskip

\textbf{Thermal Physics.}
Diffusion equations model heat transport and approach to thermal
equilibrium.

\medskip

\textbf{Continuum Mechanics.}
Tensor fields, conservation laws, and constitutive equations describe
deformable solids and fluids.

\medskip

\textbf{Relativity.}
Geometry, tensors, and invariance principles describe spacetime and
relativistic motion.

\medskip

\textbf{Quantum Mechanics.}
Hilbert spaces, linear operators, eigenvalue problems, PDEs, and probability
describe quantum states and observables.

\medskip

\textbf{Computational Physics.}
Numerical integration, PDE solvers, eigenvalue methods, and simulation make
complex physical models computationally accessible.

\end{applicationbox}

%%%%%%%%%%%%%%%%%%%%%%%%%%%%%%%%%%%%%%%%%%%%%%%%%%%%%%%%%%%%
\subsection{Exercises}
\label{subsec:mathematical-physics-exercises}
%%%%%%%%%%%%%%%%%%%%%%%%%%%%%%%%%%%%%%%%%%%%%%%%%%%%%%%%%%%%

\begin{exercise}[1]
Define velocity and acceleration from position.
\end{exercise}

\begin{exercise}[1]
State Newton's second law.
\end{exercise}

\begin{exercise}[1]
Define linear momentum.
\end{exercise}

\begin{exercise}[1]
Define kinetic energy.
\end{exercise}

\begin{exercise}[1]
Define a conservative force.
\end{exercise}

\begin{exercise}[1]
Define angular momentum.
\end{exercise}

\begin{exercise}[1]
Write the simple harmonic oscillator equation.
\end{exercise}

\begin{exercise}[1]
Define the Lagrangian.
\end{exercise}

\begin{exercise}[1]
State the Euler--Lagrange equation.
\end{exercise}

\begin{exercise}[1]
State Hamilton's equations.
\end{exercise}

\begin{exercise}[2]
Solve

\[
m\ddot{x}=F_0
\]

for constant \(F_0\), given \(x(0)=x_0\) and \(\dot{x}(0)=v_0\).
\end{exercise}

\begin{exercise}[2]
Derive the work-energy theorem in one dimension.
\end{exercise}

\begin{exercise}[2]
Find the potential corresponding to

\[
F(x)=-kx.
\]
\end{exercise}

\begin{exercise}[2]
Show that

\[
V(x)=\frac12kx^2
\]

produces Hooke's law.
\end{exercise}

\begin{exercise}[3]
Show that mechanical energy is conserved for

\[
m\ddot{x}
=
-\frac{dV}{dx}.
\]
\end{exercise}

\begin{exercise}[2]
Compute angular momentum for a particle with known
\(\mathbf{r}\) and \(\mathbf{p}\).
\end{exercise}

\begin{exercise}[3]
Show that a central force produces zero torque.
\end{exercise}

\begin{exercise}[3]
Deduce conservation of angular momentum for central-force motion.
\end{exercise}

\begin{exercise}[2]
Find the natural frequency of a mass \(m\) attached to a spring of constant
\(k\).
\end{exercise}

\begin{exercise}[2]
Find the oscillator period.
\end{exercise}

\begin{exercise}[3]
Solve the undamped oscillator with specified initial position and velocity.
\end{exercise}

\begin{exercise}[3]
Classify a damped oscillator as underdamped, critically damped, or
overdamped.
\end{exercise}

\begin{exercise}[3]
Explain resonance conceptually.
\end{exercise}

\begin{exercise}[3]
Construct the matrix equation for two coupled oscillators.
\end{exercise}

\begin{exercise}[3]
Explain why normal modes lead to an eigenvalue problem.
\end{exercise}

\begin{exercise}[2]
Construct the Lagrangian for a free particle.
\end{exercise}

\begin{exercise}[2]
Construct the Lagrangian for a harmonic oscillator.
\end{exercise}

\begin{exercise}[3]
Use Euler--Lagrange equations to recover Newton's equation for a particle in
potential \(V(x)\).
\end{exercise}

\begin{exercise}[3]
Define generalized momentum for a given Lagrangian.
\end{exercise}

\begin{exercise}[3]
Explain why a cyclic coordinate produces a conserved momentum.
\end{exercise}

\begin{exercise}[2]
Construct the Hamiltonian for

\[
L
=
\frac12m\dot{x}^2-V(x).
\]
\end{exercise}

\begin{exercise}[3]
Use Hamilton's equations to recover

\[
m\ddot{x}
=
-V'(x).
\]
\end{exercise}

\begin{exercise}[2]
Define phase space.
\end{exercise}

\begin{exercise}[3]
Compute

\[
\{q_i,p_j\}.
\]
\end{exercise}

\begin{exercise}[3]
Explain Noether's theorem conceptually.
\end{exercise}

\begin{exercise}[3]
Match time translation, spatial translation, and rotation with their
associated conservation laws.
\end{exercise}

\begin{exercise}[2]
Compute the gradient of

\[
\phi(x,y,z)
=
x^2+y^2+z^2.
\]
\end{exercise}

\begin{exercise}[2]
Compute the divergence of

\[
\mathbf{F}
=
(x,y,z).
\]
\end{exercise}

\begin{exercise}[2]
Compute the curl of

\[
\mathbf{F}
=
(-y,x,0).
\]
\end{exercise}

\begin{exercise}[2]
Compute the Laplacian of

\[
u=x^2+y^2+z^2.
\]
\end{exercise}

\begin{exercise}[3]
Explain the continuity equation

\[
\rho_t+\nabla\cdot\mathbf{J}=0.
\]
\end{exercise}

\begin{exercise}[2]
Write the one-dimensional wave equation.
\end{exercise}

\begin{exercise}[3]
Verify that

\[
u(x,t)=f(x-ct)
\]

solves the one-dimensional wave equation.
\end{exercise}

\begin{exercise}[3]
Find standing-wave modes for a string fixed at \(x=0\) and \(x=L\).
\end{exercise}

\begin{exercise}[3]
Find the corresponding normal-mode frequencies.
\end{exercise}

\begin{exercise}[2]
Write the heat equation.
\end{exercise}

\begin{exercise}[2]
Write Laplace's equation.
\end{exercise}

\begin{exercise}[2]
Write Poisson's equation.
\end{exercise}

\begin{exercise}[2]
Distinguish Dirichlet and Neumann boundary conditions.
\end{exercise}

\begin{exercise}[3]
Explain why the wave equation requires two initial conditions in time.
\end{exercise}

\begin{exercise}[3]
Explain how Fourier eigenfunctions arise in boundary-value problems.
\end{exercise}

\begin{exercise}[3]
Explain the Green's-function interpretation as point-source response.
\end{exercise}

\begin{exercise}[2]
Show that electrostatic potential satisfies

\[
\nabla^2\phi
=
-\frac{\rho}{\varepsilon_0}.
\]
\end{exercise}

\begin{exercise}[2]
State Maxwell's equations in differential form.
\end{exercise}

\begin{exercise}[3]
Explain conceptually how Maxwell's equations imply electromagnetic waves.
\end{exercise}

\begin{exercise}[3]
Derive the vacuum electromagnetic wave speed

\[
c
=
\frac{1}{\sqrt{\mu_0\varepsilon_0}}.
\]
\end{exercise}

\begin{exercise}[2]
Define a Lagrangian density conceptually.
\end{exercise}

\begin{exercise}[3]
State the field Euler--Lagrange equation.
\end{exercise}

\begin{exercise}[2]
Define the Lorentz factor

\[
\gamma.
\]
\end{exercise}

\begin{exercise}[2]
Write the Lorentz transformation along the \(x\)-direction.
\end{exercise}

\begin{exercise}[3]
Explain time dilation.
\end{exercise}

\begin{exercise}[3]
Show that

\[
\gamma\to1
\]

as

\[
v/c\to0.
\]
\end{exercise}

\begin{exercise}[2]
State relativistic momentum.
\end{exercise}

\begin{exercise}[2]
State the energy-momentum relation.
\end{exercise}

\begin{exercise}[3]
Recover

\[
E=mc^2
\]

for a particle at rest.
\end{exercise}

\begin{exercise}[3]
Explain general relativity as a geometric theory conceptually.
\end{exercise}

\begin{exercise}[2]
Define a quantum wavefunction conceptually.
\end{exercise}

\begin{exercise}[2]
State the wavefunction normalization condition.
\end{exercise}

\begin{exercise}[2]
Write the one-dimensional time-dependent Schrödinger equation.
\end{exercise}

\begin{exercise}[3]
Derive the time-independent Schrödinger equation using a separated
stationary state.
\end{exercise}

\begin{exercise}[2]
Write the position operator.
\end{exercise}

\begin{exercise}[2]
Write the one-dimensional momentum operator.
\end{exercise}

\begin{exercise}[3]
Compute the commutator

\[
[\hat{x},\hat{p}]
\]

acting on a sufficiently smooth test wavefunction.
\end{exercise}

\begin{exercise}[2]
State the position-momentum uncertainty relation.
\end{exercise}

\begin{exercise}[3]
Explain why quantum uncertainty is not simply experimental error.
\end{exercise}

\begin{exercise}[3]
Derive the eigenfunctions of the infinite square well.
\end{exercise}

\begin{exercise}[3]
Derive the infinite-well energy levels.
\end{exercise}

\begin{exercise}[2]
Find the ground-state energy for a well of width \(L\).
\end{exercise}

\begin{exercise}[3]
Explain quantum superposition.
\end{exercise}

\begin{exercise}[3]
Explain why Hilbert-space methods are natural in quantum mechanics.
\end{exercise}

\begin{exercise}[2]
State the quantum harmonic-oscillator energy levels.
\end{exercise}

\begin{exercise}[3]
Explain zero-point energy.
\end{exercise}

\begin{exercise}[2]
Use dimensional analysis to derive the scaling

\[
T\propto\sqrt{\frac{L}{g}}
\]

for a pendulum.
\end{exercise}

\begin{exercise}[3]
Explain why dimensional analysis cannot determine the factor \(2\pi\).
\end{exercise}

\begin{exercise}[3]
Nondimensionalize the damped oscillator and derive the damping ratio.
\end{exercise}

\begin{exercise}[3]
Explain why dimensionless numbers are useful for identifying physical
regimes.
\end{exercise}

\begin{exercise}[2]
Define a regular perturbation expansion.
\end{exercise}

\begin{exercise}[3]
Explain what makes a perturbation singular.
\end{exercise}

\begin{exercise}[3]
Define asymptotic equivalence

\[
f(x)\sim g(x).
\]
\end{exercise}

\begin{exercise}[3]
Explain why variational formulations are useful in physics.
\end{exercise}

\begin{exercise}[3]
Explain weak formulations conceptually.
\end{exercise}

\begin{exercise}[3]
Explain the difference between model error and numerical error.
\end{exercise}

\begin{exercise}[3]
Explain why a numerical method should sometimes preserve physical
invariants.
\end{exercise}

\begin{exercise}[4]
Analyze motion in a central potential using an effective potential.
\end{exercise}

\begin{exercise}[4]
Derive orbital equations for an inverse-square force.
\end{exercise}

\begin{exercise}[4]
Study the Kepler problem using conservation laws.
\end{exercise}

\begin{exercise}[4]
Derive normal modes of a three-mass spring system.
\end{exercise}

\begin{exercise}[4]
Study resonance of a damped forced oscillator.
\end{exercise}

\begin{exercise}[4]
Derive Euler--Lagrange equations from the stationary-action principle.
\end{exercise}

\begin{exercise}[4]
Study Hamiltonian phase-space flow.
\end{exercise}

\begin{exercise}[4]
Study Poisson brackets and canonical transformations.
\end{exercise}

\begin{exercise}[4]
Study Noether's theorem in classical mechanics.
\end{exercise}

\begin{exercise}[4]
Solve the heat equation on a finite interval by separation of variables.
\end{exercise}

\begin{exercise}[4]
Solve the wave equation with fixed-end boundary conditions.
\end{exercise}

\begin{exercise}[4]
Solve a Laplace boundary-value problem.
\end{exercise}

\begin{exercise}[4]
Construct a Green's function for a simple one-dimensional differential
operator.
\end{exercise}

\begin{exercise}[4]
Derive the electromagnetic wave equation from Maxwell's equations.
\end{exercise}

\begin{exercise}[4]
Study variational formulations of classical field theory.
\end{exercise}

\begin{exercise}[4]
Verify invariance of the spacetime interval under a Lorentz transformation.
\end{exercise}

\begin{exercise}[4]
Study four-vectors and relativistic energy-momentum.
\end{exercise}

\begin{exercise}[4]
Study geodesics in curved spacetime conceptually.
\end{exercise}

\begin{exercise}[4]
Solve the infinite square well in detail.
\end{exercise}

\begin{exercise}[4]
Study the quantum harmonic oscillator using ladder operators.
\end{exercise}

\begin{exercise}[4]
Study spectral theory of quantum observables.
\end{exercise}

\begin{exercise}[4]
Study perturbation theory in quantum mechanics.
\end{exercise}

\begin{exercise}[4]
Develop a numerical simulation of a classical oscillator and compare
energy conservation across integration methods.
\end{exercise}

%%%%%%%%%%%%%%%%%%%%%%%%%%%%%%%%%%%%%%%%%%%%%%%%%%%%%%%%%%%%
\subsection{Conceptual Summary}
\label{subsec:mathematical-physics-summary}
%%%%%%%%%%%%%%%%%%%%%%%%%%%%%%%%%%%%%%%%%%%%%%%%%%%%%%%%%%%%

Mathematical physics translates physical principles into mathematical
structures.

Classical particle mechanics begins with

\[
\boxed{
m\ddot{\mathbf{r}}
=
\mathbf{F}.
}
\]

Conservative systems can be described through potential energy:

\[
\boxed{
\mathbf{F}
=
-\nabla V.
}
\]

The same dynamics can be reformulated through the Lagrangian

\[
\boxed{
L=T-V
}
\]

and the Euler--Lagrange equations

\[
\boxed{
\frac{d}{dt}
\left(
\frac{\partial L}{\partial\dot q_i}
\right)
-
\frac{\partial L}{\partial q_i}
=
0.
}
\]

Hamiltonian mechanics uses phase space:

\[
\boxed{
\dot q_i
=
\frac{\partial H}{\partial p_i},
\qquad
\dot p_i
=
-\frac{\partial H}{\partial q_i}.
}
\]

Continuous physical systems produce PDEs such as

\[
\boxed{
u_{tt}
=
c^2\nabla^2u
}
\]

for waves,

\[
\boxed{
u_t
=
\alpha\nabla^2u
}
\]

for heat and diffusion, and

\[
\boxed{
\nabla^2u=f
}
\]

for potential problems.

Electromagnetism is organized by Maxwell's equations.

Relativity replaces absolute space and time with spacetime geometry.

Quantum mechanics replaces classical phase-space states with Hilbert-space
states satisfying

\[
\boxed{
i\hbar\psi_t
=
\hat H\psi.
}
\]

Across these theories, common mathematical themes include:

\[
\boxed{
\text{symmetry},
}
\]

\[
\boxed{
\text{conservation},
}
\]

\[
\boxed{
\text{eigenvalue problems},
}
\]

\[
\boxed{
\text{variational principles},
}
\]

\[
\boxed{
\text{scaling},
}
\]

and

\[
\boxed{
\text{approximation}.
}
\]

Mathematical physics therefore provides a major bridge between abstract
mathematics and quantitative descriptions of the physical world.

%%%%%%%%%%%%%%%%%%%%%%%%%%%%%%%%%%%%%%%%%%%%%%%%%%%%%%%%%%%%
\subsection{Section Connections}
\label{subsec:mathematical-physics-connections}
%%%%%%%%%%%%%%%%%%%%%%%%%%%%%%%%%%%%%%%%%%%%%%%%%%%%%%%%%%%%

\chapterconnection{
Mathematical Physics applies the modeling principles of
Section~\ref{sec:mathematical-modeling} to physical laws. Calculus and
differential equations describe motion and fields, linear algebra produces
normal modes and quantum eigenstates, vector calculus expresses flux and
field laws, variational methods generate equations of motion, geometry
provides the language of relativity, and functional analysis supports
quantum mechanics. The next section, Continuum Mechanics, develops the
mathematics of continuously distributed matter through deformation,
kinematics, stress, strain, constitutive laws, and conservation of mass,
momentum, and energy.
}

%%%%%%%%%%%%%%%%%%%%%%%%%%%%%%%%%%%%%%%%%%%%%%%%%%%%%%%%%%%%
% SECTION SOURCE TRACKING
%%%%%%%%%%%%%%%%%%%%%%%%%%%%%%%%%%%%%%%%%%%%%%%%%%%%%%%%%%%%

%------------------------------------------------------------
% 13.3 Mathematical Physics
%------------------------------------------------------------
%
% Source basis:
%   - Standard classical mechanics.
%   - Standard analytical mechanics.
%   - Standard wave, heat, and potential theory.
%   - Standard electromagnetic theory.
%   - Introductory special and general relativity.
%   - Introductory quantum mechanics.
%   - Standard variational, asymptotic, and numerical methods in physics.
%
% Used for:
%   - physical quantities and units;
%   - kinematics;
%   - Newtonian mechanics;
%   - momentum;
%   - work;
%   - kinetic and potential energy;
%   - conservative forces;
%   - energy conservation;
%   - power;
%   - angular momentum;
%   - central forces;
%   - inverse-square forces;
%   - harmonic oscillators;
%   - damping;
%   - forcing;
%   - resonance;
%   - coupled oscillators;
%   - normal modes;
%   - Lagrangian mechanics;
%   - action principles;
%   - Euler--Lagrange equations;
%   - generalized momentum;
%   - cyclic coordinates;
%   - Hamiltonian mechanics;
%   - Hamilton's equations;
%   - phase space;
%   - Poisson brackets;
%   - symmetry;
%   - Noether's theorem;
%   - scalar and vector fields;
%   - gradient, divergence, curl, and Laplacian;
%   - continuity equations;
%   - wave equations;
%   - standing waves;
%   - eigenfunction expansions;
%   - heat equations;
%   - Laplace and Poisson equations;
%   - boundary and initial conditions;
%   - Fourier methods;
%   - Green's functions;
%   - electrostatics;
%   - Maxwell's equations;
%   - electromagnetic waves;
%   - field theory;
%   - field Euler--Lagrange equations;
%   - stress and strain preview;
%   - special relativity;
%   - Lorentz transformations;
%   - spacetime intervals;
%   - relativistic momentum and energy;
%   - general relativity preview;
%   - Einstein field equations preview;
%   - quantum states and wavefunctions;
%   - Schrödinger equation;
%   - operators and observables;
%   - expectation values;
%   - commutators;
%   - uncertainty relations;
%   - infinite square well;
%   - quantum harmonic oscillator;
%   - Hilbert spaces;
%   - dimensional analysis;
%   - nondimensionalization;
%   - dimensionless parameters;
%   - perturbation theory;
%   - singular perturbations;
%   - asymptotic methods;
%   - variational methods;
%   - weak formulations;
%   - numerical mathematical physics;
%   - model and numerical error;
%   - applications and exercises.
%
% Notes:
%   - Continuum Mechanics is developed next in Section 13.4.
%   - Fluid Dynamics follows in Section 13.5.
%   - Control Theory follows in Section 13.6.
%   - Network Science follows in Section 13.7.
%
%%%%%%%%%%%%%%%%%%%%%%%%%%%%%%%%%%%%%%%%%%%%%%%%%%%%%%%%%%%%